\documentclass[12pt]{article}

% House style preamble
\input{.house-style/preamble.tex}

\hypersetup{
  pdftitle={Access transparency in grammar: Mediated accessibility as a comparative concept},
  pdfkeywords={access transparency, mediated accessibility, agreement, comparative concepts, number transparency, transparent free relatives, projectibility}
}

\title{Access transparency in grammar: Mediated accessibility as a comparative concept}
\author{Brett Reynolds \orcidlink{0000-0003-0073-7195}%
\thanks{Contact: \href{mailto:brett.reynolds@humber.ca}{brett.reynolds@humber.ca}}\\
Humber Polytechnic \& University of Toronto
\and
Jong-Bok Kim\\
Kyung Hee University}
\date{\today}

% Tables: longtable + booktabs + calc + array (used by section .tex files).
\usepackage{longtable}
\usepackage{booktabs}
\usepackage{calc}
\usepackage{array}
\usepackage[modulo,mathlines]{lineno}
% Pandoc-style longtables set \LTcaptype to none; lineno expects that counter.
\newcounter{none}
% Tightlist for compact lists; pandocbounded/real are legacy compatibility
% from earlier markdown→pandoc pipeline; harmless if unused.
\providecommand{\tightlist}{%
  \setlength{\itemsep}{0pt}\setlength{\parskip}{0pt}}
\providecommand{\real}[1]{#1}
\providecommand{\pandocbounded}[1]{#1}

\definecolor{schemablue}{RGB}{239,247,255}
\definecolor{schemabluerule}{RGB}{185,211,230}
\newsavebox{\schemaboxbox}
\newenvironment{schemabox}{%
  \begin{center}
  \setlength{\fboxsep}{8pt}%
  \setlength{\fboxrule}{0.4pt}%
  \begin{lrbox}{\schemaboxbox}%
  \begin{minipage}{0.88\linewidth}
  \centering
}{%
  \end{minipage}%
  \end{lrbox}%
  \fcolorbox{schemabluerule}{schemablue}{\usebox{\schemaboxbox}}%
  \end{center}
}

% \Cref / \cref auto-prefix section references with the right word/symbol.
\usepackage[capitalize,nameinlink]{cleveref}
\crefname{section}{§}{§§}
\Crefname{section}{§}{§§}

% Cosmetic fixes for the build.
\setlength{\headheight}{26pt}        % silence fancyhdr headheight warning
\setlength{\emergencystretch}{3em}   % allow stretchier line-breaks in bibliography notes

% Prevent line breaks inside inline math at relations and binary operators.
\relpenalty=10000
\binoppenalty=10000

\begin{document}
\maketitle
\linenumbers

\begin{abstract}
Linguists use \term{transparency} for cases in which a relation can access a property through material that could have hidden it. This broad pattern is \term{mediated accessibility}; the paper defines its narrower grammatical relation, \term{access transparency}, as a comparative concept. An application passes four gates: an independently diagnosable bearer, a distinct mediator, evidence that the relation tracks the property, and an independently established concealing counterpart in the same constructional family. Verdicts are positive, non-instance, or indeterminate, and feature summation, whole-configuration construal, and external bearers come out as non-instances. English \term{number-transparent quantificational nouns} (QNs), whose agreement can track the \mention{of}-complement rather than the first noun, supply the main empirical case, with a corpus pilot giving partial prospective support and a feasibility limit; transparent free relatives supply a second illustration, positive under Kim's constructional analysis and indeterminate analysis-neutrally. Spanish and Tsez extend the protocol cross-linguistically. \term{Projectibility} is the warrant for extending observed regularities to specified new cases, but it's secondary here: the comparative concept predicts nothing by itself, and a language-particular covariance pattern supports that extension only when an observable base warrants non-trivial inferences. An identifiable \term{stabilizer} is a mechanism that would separately explain why the covariance holds. The result distinguishes a relational analysis that earns its keep from a label that restates the facts it describes.

\end{abstract}

\noindent\textbf{Keywords:} access transparency; mediated accessibility; agreement; comparative concepts; number transparency; transparent free relatives; projectibility

\section{Introduction}\label{sec:1}

Two simple sentences make the puzzle visible. \mention{A lot of tourists are in town.} \mention{A lot of money is being spent.} The verb agrees with the \mention{of}-complement (\mention{tourists}, \mention{money}) rather than with the surface head \mention{lot}, which is morphologically singular in both cases. The complement is plural count in the first sentence, singular non-count in the second. The construction lets the relevant feature of the complement reach the verb.

\term{Number-transparent quantificational nouns} is the term used in Huddleston and Pullum's \mention{Cambridge Grammar of the English Language} \autocite*{huddleston-pullum-2002} (henceforth CGEL). Core members include \mention{lot, plenty, lots, bags, heaps, loads, oodles, stacks, remainder, rest, number,} and \mention{couple}, with \mention{majority} and \mention{minority} at the boundary. Membership is small and relatively stable, and members head NPs whose agreement override is, on CGEL's analysis, the default pattern, contrasting with the optional override of ordinary collective nouns. Many other configurations behave similarly: a form mediates between an internal property and a grammatical or interpretive relation, and the property remains accessible through the mediator. The umbrella schema for the family follows.

\begin{quote}
$X$ is \mention{transparent} with respect to property $P$ for relation $R$ when $P$ remains accessible through $X$ for the purposes of $R$.
\end{quote}

The schema is a tool for a particular question: when is calling something \mention{transparent} doing analytical work, and when is it just a label? Sometimes the label genuinely earns its keep; sometimes it doesn't. The \mention{a lot of} construction shows what earning the label looks like.

The schema's substantive content is the \mention{remains accessible through $X$} clause. A case fits only if $X$ has internal structure whose default behaviour would have made $P$ inaccessible to $R$. Cases without that default-conceal background are excluded. Coordination (§\ref{sec:2-3}) is the textbook exclusion: \mention{Bob and Jane are happy} has no head whose morphological number competes for agreement with the conjuncts' features, so there's nothing the construction could have hidden, and calling coordination \mention{transparent} would be a relabel rather than an analysis. Everyday senses of \term{transparent prose} or \term{transparent intentions}, where no specific mediator-property-relation triple is in play, are excluded for a related reason: the schema's slots have no values to bind. The exclusion mechanism is what keeps the schema from cataloguing every use of \mention{transparent} in linguistics, and the rest of the paper is occupied with cases where the slots have non-trivial values.

Once we recognise that \mention{lot, plenty, lots, bags, heaps, loads, oodles, stacks, remainder, rest, number,} and \mention{couple} (in its quantificational sense) form a small, conventionally stable class, we can predict more than the agreement override that motivated the label. Members of the class show characteristic patterns of determiner co-occurrence; they admit partitive \mention{of}-complements in some subgroups but not others; they behave distinctively under the fused-head construction; the count-or-non-count profile of the \mention{of}-complement determines the profile of the whole NP; modification by ordinary adjectives is restricted in characteristic ways. Knowing that a noun belongs to the class lets us predict much of the rest, with the predictions reliable for the core members and graded at the boundary (\mention{majority, minority}, sense-split \mention{couple}). The label is doing work: a small lexical inventory whose membership predicts a cluster of correlated behaviours, attenuated at the edges rather than guaranteed everywhere.

Cases like this we'll call \term{projectibility profiles}. A profile has three components: an \term{observable input} (here, lexical-class membership); a set of \mention{predictions} the input licenses (here, the cluster of agreement, distributional, and modificational behaviours); and a \term{stabilising process} that holds the input together (here, ordinary lexical inheritance). The number-transparent class is a strong example. The predictions are non-trivial; the inventory is stable; the stabilising process is the same one speakers use for any closed grammatical class.

Not every transparency label names a profile this strong. Some name profiles whose predictions are weaker, or whose stabilising process is different, or both. The paper asks which transparency labels do projection-supporting work, and what makes them able to do it.

$X$ (the mediator) can be a lexical item, a partly or fully schematic construction, an intensional context, or a derivationally complex word. $P$ can be a number/count feature, a constituent's meaning, a referent, or a stem's lexical representation. $R$ can be agreement, interpretation, substitution under co-reference, or masked priming. Accessibility itself admits two readings: class-like (core members behave categorically, even where class boundaries are fuzzy) and gradient (compositionality scales, transparency ratings, priming magnitudes).

The argument is constructive. The paper isn't a survey of every use of \term{transparency} in linguistics, nor a paper on extraction, islands, barriers, or phases; extraction-transparency is acknowledged in §\ref{sec:2} but not developed. It isn't a paper on phonological or orthographic transparency, nor on the everyday senses of \mention{transparent explanation} or \mention{transparent terminology}. Referential or \term{de re}/\term{de dicto} transparency receives a controlled mention in §\ref{sec:4} as the boundary case where the schema's mediator-form structure stretches thinnest.

\term{Mediated accessibility} is more specific than \term{clarity} and broader than \term{compositionality}. Clarity is rhetorical: it doesn't single out a mediator or a relation. Compositionality is the limiting case of \term{form-meaning} transparency, where a mediator's parts predict the whole. The umbrella schema is broader: it covers form-meaning cases (constructions, compounds, derived words) and accessibility cases that aren't reducible to whether parts predict wholes. Number transparency is the paradigm of the second type: \mention{lot}'s morphological singular doesn't predict the verb's number, and under the construction it doesn't have to.

Section 2 takes up morphosyntactic transparency, with number-transparent quantificational nouns and transparent free relatives as the paradigm cases. Sections 3 to 5 work through the form-meaning accessibility family: theoretical compositionality \autocite{goldberg-2006-constructions-at-work}, rated compound transparency \autocite{bell-schafer-2016-modelling-transparency}, and time-sensitive lexical processing \autocite{heyer-kornishova-2018-priming-eventually}. Section 6 applies Haspelmath's \autocite*{haspelmath-2010-comparative-concepts} comparative-concept discipline to the schema's cross-linguistic ambitions. Section 7 names the stabilising structures behind each profile and treats Slater's SPC framework \autocite{slater-2015-natural-kindness} and Khalidi's causal-network account \autocite{khalidi-2018-causal-networks} as complementary diagnostics for projection-supporting stability.

\section{Morphosyntactic transparency}\label{sec:2}

Morphosyntactic transparency is the case where $X$ is a syntactic construction (or a lexical class within one) and $R$ is a grammatical relation sensitive to a property borne inside $X$ but not by its head. Feature matching between a non-head property and a co-occurring exponent is the paradigm, and agreement is the paradigm of that. Selection, category licensing, binding, and dependency formation are possible targets, as §\ref{sec:2-5} shows for transparent free relatives.\footnote{Extraction transparency (whether an island or barrier is opaque or transparent to long-distance dependency formation) is a morphosyntactic-accessibility relation under the schema (with $R$ = dependency formation), and the associated question remains active. Transformational accounts posit empty elements (traces, copies) that mediate the connection between a fronted element and its base position; non-transformational dependency accounts \autocite{gibson2025,dacunha-gibson-2025-syntactic-complexity} argue that no such mediating empty elements are needed and that direct head-to-dependent connections better fit acceptability and processing data. Both sides make schema-level claims about whether $X$ contains invisible mediating structure. The centre of the argument here is the agreement-style cases.} Subject-verb agreement is the textbook instance. The concrete contrasts are simple. In (\ref{ex:cgel-two-direction}a), agreement tracks \mention{spots} through singular \mention{number}; in (\ref{ex:tfr-r-relations}a), agreement tracks \mention{trousers} through singular \mention{what}.

This section carries the paper's primary field-relative claim. The projected targets are syntactic and morphosyntactic: agreement, syntactic-category licensing, complement selection, behaviour under complement ellipsis, and distributional restrictions inside NP and fused-relative constructions. Nothing in the syntactic result depends on phrase-structure representation as such: the profiles require only observable relations among a mediator, an internal property, and the grammatical relation that tracks it, whether those relations are represented in constituency or dependency terms. The section contains two grammatical cases, but only the QN section carries a prospective candidate-member test; the TFR result remains a conditional illustration based on Kim's constructional analysis and pilot corpus evidence.

\subsection{Number-transparent quantificational nouns}\label{sec:2-1}

The paradigm again is \textit{CGEL}'s two-direction contrast in (\ref{ex:cgel-two-direction}): plural agreement despite singular \mention{number}, and singular agreement despite plural \mention{heaps}. Schema-wise (under the umbrella schema), $X$ is the number-transparent quantificational noun + \mention{of}-complement construction; $P$ is the number/count profile of the \mention{of}-complement; $R$ is subject-verb agreement.

\textit{CGEL}'s \enquote{simple agreement rule} is head-driven: \enquote{the verb agrees with a subject with the form of an NP whose person--number classification derives from its head noun} (§18.1, p.~499), as in (\ref{ex:simple-agreement}).

\ea\label{ex:simple-agreement}
\ea \mention{The nurse wants to see him.}
\ex \mention{The nurses want to see him.}
\z
\z

Number-transparent NPs override the rule. \textit{CGEL} describes the override as obligatory in the central cases; rare singular-verb attestations with \mention{number} are noted in footnote 73 to §18.2 (p.~502) and analyzed there as hypercorrections rather than as variant grammar. The contrast with the optional override of collective nouns (§\ref{sec:2-2}) is the operative point. The construction's job is to let the \mention{of}-complement's number/count profile reach agreement past the head's morphological number, which would otherwise have determined it.

The core inventory is small and well documented. \textit{CGEL} identifies the principal members in example [57] (§3.3, p.~350): \mention{lot, plenty} (singular forms; \mention{lot} most often with \mention{a}-determiner, \mention{plenty} typically without a determiner); \mention{lots, bags, heaps, loads, oodles, stacks} (plural forms, typically without a determiner); \mention{remainder, rest} (singular forms typically with \mention{the}-determiner); \mention{number, couple}₂ (singular forms typically with \mention{a}-determiner, taking only plural \mention{of}-complements, partitive or pseudo-partitive). Of these, \enquote{the clear cases of number-transparent singular nouns are \mention{lot, number}, and \mention{couple}₂} (p.~503); \mention{majority} and \mention{minority} are borderline: plural override is \enquote{surely obligatory} in examples like \mention{The majority of her friends are Irish}, but singular agreement is occasionally attested (§18.2, pp.~503--504). The determiner patterns above are typical, not absolute: \mention{the lot of it} (in a totality sense) is grammatical alongside \mention{a lot of it}, and similar sense-conditioned variations exist for other members.

Keizer's \autocite*{keizer-2007-english-noun-phrase} study of English NP categorization and Van Eynde and Kim's \autocite*{van-eynde-kim-2023-pseudo-partitives} Head-driven Phrase Structure Grammar (HPSG) analysis of pseudo-partitives are the immediate background for treating these strings as structured binominal configurations rather than as a single loose family. Within the binominal family, Kim \autocite*[98--103]{kim-2026-form-function} distinguishes three sub-types by agreement head, illustrated in (\ref{ex:three-subtypes}): measure nouns (\mention{pound}), where the first noun controls agreement; quantity nouns (\mention{lot}, \mention{number}), where the second noun controls it; and collective nouns (\mention{group}, \mention{bunch}), which allow either.

\ea\label{ex:three-subtypes}
\ea Measure (N1 controls): \mention{One pound of tomatoes equals 90 calories.}
\ex Quantity (N2 controls): \mention{A lot of marbles are about to fall on the floor.}
\ex Collective (either): \mention{A group of police cruisers is speeding across the bridge}; \mention{a group of immigrants move in}.
\z
\z

The three example sentences are Kim's, from COCA \autocite[99--102]{kim-2026-form-function}. The quantity type is \textit{CGEL}'s number-transparent class, and the measure type gives the schema its concealing variant: the same binominal string with agreement held by the first noun (§\ref{sec:2-3}).

Class membership predicts a cluster of behaviours beyond the agreement override. The relevant contrast is the standard \term{partitive} / \term{pseudo-partitive} distinction: partitive \mention{of}-complements require a definite complement that picks out a subset of a definite set (\mention{of the money}, \mention{of those books}, \mention{of these soldiers}), while pseudo-partitive \mention{of}-complements express a quantity, measure, or collection relation with a bare or otherwise non-definite complement in the QN frame (\mention{of money}, \mention{of soldiers}, \mention{of some paintings}) \autocite{selkirk-1977-noun-phrase-structure,stickney-2007-pseudopartitive-partitive,kim-michaelis-2020-syntactic-constructions,van-eynde-kim-2023-pseudo-partitives}.\footnote{The headedness issue is framework-internal. HPSG analyses treat the first noun as syntactic head in \mention{A of the B} partitives, while \textit{CGEL} retains the quantificational noun as head and treats agreement with the complement as override. Since the schema is framework-portable rather than neutral (§\ref{sec:3}), the paper doesn't adjudicate the headedness choice; it records that the verdict is conditional on it.}

\mention{Lot} most often takes \mention{a} as determiner with limited premodification (\mention{A whole/huge lot of time has been wasted}; \textit{CGEL}, §3.3, p.~350); \mention{the lot of it} is also available in a totality sense, but other determiners are blocked (*\mention{some lot of it}). \mention{Plenty} typically takes no determiner or modifier (\mention{plenty of butter}, not generally \mention{a remarkable plenty of butter}), and like the plural-only forms it strongly prefers the bare \mention{of}-complement.\footnote{A COCA pilot (May 2026, extended July 2026) supports the preference but not a rate. The bare complement dominates for \mention{plenty} and the plural-only forms, while \mention{rest} and \mention{remainder} prefer the determined one on matched complements ($\approx\!97\%$ for \mention{day, time, people, money}). No overall share is defensible: the two frames were sampled over different complement inventories, and frequency in any case supplies no denominator of occasions on which a partitive reading was intended. The partitive frame is attested for every member of both subgroups. Full counts and methodology in Appendix~\ref{a.2-partitive-vs.-pseudo-partitive-of-complements}; raw data in the supplementary material (see Data availability).} The plural-only forms (\mention{lots, bags, heaps, loads, oodles, stacks}, and \mention{piles} beyond \textit{CGEL}'s list) show the same pattern.

\ea\label{ex:plural-only-of-complements}
\ea \mention{They have oodles of money.}
\ex \mention{Oodles of the money had already been spent.}
\z
\z

The contrast in (\ref{ex:plural-only-of-complements}) shows the plural-only pattern: the first is the characteristic use, the second a marked one. It is not a blocked one. A wildcard sweep of the determined frame attests it for every member of the subgroup, at 465 tokens across 374 distinct complements, where a probe over four hand-picked complements had returned 13 and two zeros (Appendix~\ref{a.2-partitive-vs.-pseudo-partitive-of-complements}). \mention{Remainder} and \mention{rest} are a different case again; \mention{number} and \mention{couple}₂ take plural obliques.

\ea\label{ex:partitive-plural-oblique}
\ea \mention{the remainder of the time}
\ex \ungram{\mention{the remainder of time}} (partitive reading only; fine as \enquote{all time henceforth})
\ex \mention{the rest of society}
\ex \mention{a number of the protesters}
\ex \ungram{\mention{a number of money}}
\z
\z

The obliques in (\ref{ex:partitive-plural-oblique}) are structured, but the lexical classification of \mention{rest} is not settled by the corpus pattern. Under the semantic criterion adopted here, \mention{rest} and \mention{remainder} denote the complement subset of a presupposed whole rather than a quantity of it, so the pseudo-partitive construction, defined as one in which N1 denotes \enquote{a quantity or amount} of what N2 denotes \autocite[271]{van-eynde-kim-2023-pseudo-partitives}, is unavailable. A competing analysis could treat bare \mention{the rest of society} as pseudo-partitive by treating \mention{rest} as a part or amount expression. The data establish the bounded-and-unique condition on the bare frame, but do not decide that lexical-semantic question; Table~\ref{tab:qn-matrix} therefore marks \mention{rest} as unresolved while retaining the semantic analysis as the working description.

For the bare-complement reading, the whole must be uniquely identifiable \mention{and} bounded, so that it has a proper part to take a remainder of. Generic collectives supply both unaided, which is why (\ref{ex:partitive-plural-oblique}c) needs no determiner and is common: \mention{the rest of society} alone has 442 COCA tokens, with \mention{humanity} 248, \mention{nature} 78, \mention{mankind} 77, \mention{creation} 42, and \mention{Congress} 21 behind it. An unbounded complement supplies neither. Bare \mention{of time} offers no partition, which is why (\ref{ex:partitive-plural-oblique}b) is starred on the relevant reading and why the string survives only on the unrelated duration reading, \enquote{all time henceforth}. That reading accounts for most of the corpus tokens of \mention{the rest of time}, so those tokens are not evidence about the partitive/pseudo-partitive classification (Appendix~\ref{a.2-partitive-vs.-pseudo-partitive-of-complements}).

\textit{CGEL}'s own minimal pair for \mention{lot} has the complement controlling agreement in both directions (p.~349).

\ea\label{ex:override-visibility}
\ea \mention{A lot of work was done.}
\ex \mention{A lot of errors were made.}
\z
\z

The pair in (\ref{ex:override-visibility}) shows what number transparency \mention{is} in the umbrella schema's vocabulary. $X$ (the construction) leaves $P$ (the number/count profile of the \mention{of}-complement) accessible to $R$ (agreement). What $X$ conceals from agreement is the morphological number of the head noun \mention{lot, plenty, heaps}, etc. The agreement-targeting feature is on the \mention{of}-complement, not on the surface head. That's what makes $X$ transparent with respect to $P$: $X$ lets $P$ reach the verb anyway.

The override is only visible where the head's number and the complement's diverge. (\ref{ex:override-visibility}b) shows it (singular \mention{lot}, plural \mention{errors}); (\ref{ex:override-visibility}a) doesn't, since head and complement both call for the singular. The discriminating cases are the crossed ones: a singular-form quantificational noun with a plural complement, or a plural-form one (\mention{lots, heaps}) with a non-count complement. The matched cases are silent, not counterevidence. Transparency is therefore under-detected: a class can be transparent while most of its tokens look head-driven, so corpus frequency understates it.

\subsection{Collective NPs as a transitional case}\label{sec:2-2}

Singular collective nouns (\mention{committee, jury, family, team, government}; also \mention{clergy, intelligentsia, public}, etc.) allow optional plural override of head-driven agreement, as in \textit{CGEL} [8] and [10] (§18.2, pp.~501--502).

\ea\label{ex:collective-agreement}
\ea \mention{The committee has}/\mention{have} \mention{not yet come to a decision.}
\ex \mention{The committee consists}/\ungram{\mention{consist}} \mention{of two academic staff and three students.}
\ex \mention{One committee, appointed last year, has}/\marg{\mention{have}} \mention{not yet met.}
\z
\z

The override in (\ref{ex:collective-agreement}) is construal-sensitive: singular agreement tracks the unit construal, plural the member construal. Plural is unavailable with predicates that apply to the whole but not individual members, and is \enquote{of questionable acceptability} if the collective takes determiner \mention{one} or \mention{another} (p.~502).

Both collective NPs and number-transparent quantificational NPs override the head's morphological singular at agreement, and the override is possible because plural agreement requires only non-singular construal, not the exact enumeration that singular \mention{a}(\mention{n}) and low cardinals demand. That is why tight determiners degrade the override that (\ref{ex:collective-agreement}a) allows while plural agreement survives it: the override operates at the loose end of the count cluster, not across it. Reynolds \autocite*{reynolds2026countability} develops that loose/tight dissociation at length.

The two cases differ in where the non-singular construal comes from. In collectives it is lexically encoded in the head, so the override is selective: it needs a predicate that licenses member-level construal, which is why whole-unit predicates block it. In number-transparent QNs it is imported from the \mention{of}-complement, so the override applies regardless of predicate type. That difference matters more than it first appears, because only one of them satisfies the schema as stated. The schema requires one element within $X$ to bear $P$ and another to conceal it. In the QN case those are distinct expressions: the \mention{of}-complement bears the number profile, the quantificational head conceals it. In \mention{the committee have decided} they are not. The singular morphology and the member-level plurality are two representations of one lexical head, so there is no second element doing the concealing. Collectives are therefore a near case: a competition between morphological and semantic properties within a single sign, not access across a mediator. Broadening \mention{element} to cover representational layers would accommodate them, at the cost of the non-instance test's grip, since almost any property could then be said to conceal another representation of itself. The transitional label in this subsection's title is meant in that strict sense.

Alongside \textit{CGEL}'s override framing, the HPSG tradition encodes the same observation as index rather than concord agreement: concord tracks morphosyntactic features internal to the NP, while index agreement can track the referential or semantic index available to the agreement target \autocite{wechsler-zlatic-2003-agreement,kim-2004-hybrid-agreement}. Kim and Sells \autocite*[64--65]{kim-sells-2015-binominal-nps} run the two levels on the cases at issue here: measure \mention{four pounds was} is morphosyntactically plural but index-singular, while collective \mention{this team have} is the mirror image, morphosyntactically singular with a plural index. On that encoding, collective agreement makes semantic plurality available to agreement, but it isn't access transparency in the strict sense used here, because the bearer and the proposed mediator aren't distinct elements. The index/concord split is consistent with the framework-portability of §\ref{sec:3}.

\textit{CGEL} itself notes that \enquote{the division between collective and number-transparent nouns is by no means sharply drawn} (§18.2, p.~502). Some nouns straddle. \mention{Couple} has two senses: \mention{couple}₁ (a man and woman in a relationship) is a regular collective; \mention{couple}₂ (`approximately two') is number-transparent. \mention{Bunch} heads the same binominal \mention{of}-frame as §\ref{sec:2-1}'s quantity nouns but is a collective noun in Kim's \autocite*[98--103]{kim-2026-form-function} classification, the sub-type where either noun can control agreement; the straddling \textit{CGEL} reports is what that classification predicts. \mention{Bunch} is also the case where the profile framework does demotion work: \textit{CGEL} claims it tilts collective for inanimate aggregates but transparent for groups of people, illustrated by the contrast in (\ref{ex:bunch-transition}).

\ea\label{ex:bunch-transition}
\ea \mention{A bunch of flowers was presented to the teacher.}
\ex \mention{A bunch of hooligans were seen leaving the premises.}
\z
\z

\textit{CGEL} treats the first as collective and the second as transparent (§18.2, p.~503). A pilot COCA check (May 2026) supports the animate half, where \mention{bunch}-subjects are consistently plural, and finds no positive support for the inanimate half: inanimate \mention{bunch} is rare in subject position, \textit{CGEL}'s own example is unattested, and the few tokens that surface take plural rather than the predicted singular. Corpus absence of a constructed example isn't ungrammaticality, so this doesn't touch \textit{CGEL}'s contrast as a claim about grammar; what it shows is that the inanimate-collective pattern doesn't project as a corpus regularity. That localizes the projectibility claim rather than the grammatical one, and marking where a textbook contrast doesn't project is what demotion conditions are for. Counts in Appendix~\ref{a.4-bunch-animate-vs.-inanimate-aggregates}.

The paradigm and the transitional case also differ in internal structure. Number-transparent NPs are class-like, with the override the default behaviour for class members. Collective NPs are gradient: the override probability depends on the head, the predicate, the determiner, the dialect register, and the construal the speaker has in mind. A theory of morphosyntactic transparency that ignored the difference between a class-based default override and a construal-driven optional one would lose the interesting structure of both cases.

\subsection{Coordination as a non-instance}\label{sec:2-3}

Coordination of singulars triggers plural agreement.

\ea\label{ex:ordinary-coordination}
\mention{Bob and Jane are happy.}
\z

The ordinary coordination in (\ref{ex:ordinary-coordination}) looks like a transparency case at first (the conjuncts' number features percolate up to the coordinate construction and determine agreement), but the schema excludes it. The schema requires that $X$ \mention{could have} made $P$ inaccessible. The coordinate construction does no such hiding: there's no head whose morphological number competes with the conjuncts' features for agreement.\footnote{This verdict presupposes a headless (or numberless-head) analysis: with no conjunct serving as a number-bearing head, nothing conceals the conjuncts' summed features. The picture changes under an analysis on which one conjunct is the head. Closest-conjunct agreement \autocite{bhatt-walkow-2013} is the case in point: where a verb tracks the nearest conjunct, that agreement conceals the resolved features and resolved agreement reaches them past it, so the same test can return an instance. As with the language-internal concealing baseline of §\ref{sec:6}, the verdict rides on the prior syntactic analysis.} The plural agreement is compositional addition (1 + 1 = 2), not mediation.

Some coordinate NPs take singular agreement.

\ea\label{ex:holistic-coordination}
\ea \mention{Bacon and eggs is my favourite breakfast.}
\ex \mention{Wine and cheese is expected.}
\ex \mention{Fish and chips is expensive these days.}
\z
\z

Some of the same coordinate strings also take plural agreement.

\ea\label{ex:coordination-plural-counterparts}
\ea \mention{Wine and cheese make the perfect party pair.}
\ex \mention{Fish and chips taste their best right after frying.}
\z
\z

The singular holistic cases in (\ref{ex:holistic-coordination}) and the plural counterparts in (\ref{ex:coordination-plural-counterparts}) leave the exclusion intact. In such cases agreement tracks a construal of the coordinate phrase: either a dish, package, event type, or conventional pairing, or a plural set of jointly mentioned items. The singular feature isn't a property of either conjunct that remains accessible through the coordinate construction; it's a construal of the coordinate phrase as a whole. Ordinary plural coordination and holistic singular coordination therefore differ in agreement semantics, but neither is transparent in the schema's sense. The former is compositional summation; the latter is whole-phrase unit construal.\footnote{Cases like \mention{One and one is two} are an arithmetic-expression construction, not a meal/package coordination: \mention{and} functions close to an addition operator, and the subject denotes an operation, equation, or expression. Not transparency either.} \textit{CGEL} treats coordination-and-agreement under a separate heading (§18.4), not alongside the override constructions of §18.2.

The clean distinction across §\ref{sec:2}'s cases: in number-transparent QNs, complement number becomes accessible through a quantificational head; in collective plural agreement, member plurality is available to agreement but is borne by the head itself, so no distinct mediator intervenes; in unit-construal coordination, the whole coordinate phrase is construed as singular and no internal property is being made transparent. Only the first is access through a mediator.

The \enquote{$X$ could have made $P$ inaccessible} clause works as an independent diagnostic, not an ad hoc filter. The test asks whether the construction-type admits a concealing counterpart: is there a version of the same construction-type, attested in the language, in which $X$'s internal structure makes $P$ inaccessible to $R$? For English NPs, head-driven agreement is the default and the number-transparent class is the marked exception within the same NP type. For coordinate constructions no such counterpart exists: the type admits no version in which conjunct features fail to reach agreement, so coordination has no control against which access can be measured. The diagnostic distinguishes mediator-types from compositional-summation types: access is informative only where the same mediator type has opaque or partially opaque alternatives.

One condition has to be attached, or \mention{same construction-type} becomes adjustable after the fact. Drawn narrowly enough, no transparent construction has an opaque sibling and every case is a non-instance; drawn broadly enough, almost anything does and the filter passes everything. So the concealing control must belong to a constructional family that is independently motivated, established on criteria that don't appeal to the transparency effect the test is assessing. English head-driven agreement qualifies: it is the unmarked pattern for NP subjects on grounds that have nothing to do with quantificational nouns, which is what makes the number-transparent class a marked exception \mention{within} a family fixed elsewhere. Tsez class IV agreement with a clausal argument qualifies on the same footing. Where no such family can be identified without invoking the effect, the test returns no verdict rather than a negative one; the typological cases in §\ref{sec:6} show why that control has to be fixed language-internally.

The condition also locates where the diagnostic is soft. Deciding what counts as the same construction-type is itself a language-internal analysis, so the test is crisp where the family is independently motivated and loose where it isn't. That is a real limitation on the paper's central instrument, not a detail about its edges.

A second non-instance fails on a different role, and the contrast is worth having because it answers the obvious deflationary reading of the QN case. In deferred or metonymic reference, agreement can track a referent the sentence never names.

\ea\label{ex:reference-transfer}
\ea \mention{The hash browns at table nine are}/\ungram{\mention{is}} \mention{getting cold.}
\ex \mention{The hash browns at table nine is}/\ungram{\mention{are}} \mention{getting angry.}
\z
\z

Kim's \autocite*[1108]{kim-2004-hybrid-agreement} pair, which he credits to Nunberg, turns on which entity the subject is used to pick out: the food in the first, the customer who ordered it in the second. Warren \autocite*[16]{warren2006} states the generalization, that metonymic subjects \enquote{need not agree as to number with their predicates} where metaphorical ones consistently do, and adds that a metonymic pronoun \enquote{sometimes agrees with the explicit and sometimes with the implicit element of the expression}, the choice settled by context.

Filled into the frame of §\ref{sec:1}, the \textit{bearer} cell is empty. The singular in (\ref{ex:reference-transfer}b) is sourced from the designated customer, who is expressed nowhere in the subject NP, so there is no element within $X$ bearing $P$ and nothing for the relation to reach \mention{through}. Coordination fails the test for want of a mediator; this fails for want of a bearer. Both are outside the schema, for different reasons.

Number transparency can be read as construal-tracking: the verb agrees with how the speaker conceives the subject, rather than with any constituent. Reference transfer shows what construal-tracking looks like on its own, and it looks different. It's context-dependent and reversible on one string, it isn't tied to a lexical class, and on Warren's account the pronoun may go either way. The QN override is none of these: near-categorical for core members, tied to a small conventional class, and construction-linked, as Appendix~\ref{sec:A-5}'s \mention{a number of} versus \mention{the number of} control shows on a single lexeme. Same surface phenomenon, different mechanism, and the frame's cells say which.

\subsection{Number-transparent NPs as a projectibility profile}\label{sec:2-4}

Section 1 set the secondary question: does identifying a positive access relation let us predict behaviour? For number-transparent quantificational nouns the profile hypothesis can be tested in two task settings: description of known members and class discovery for candidates. (The framework returns negatives too: the inanimate half of the \mention{bunch} contrast fails to project, §\ref{sec:2-2}, and coordination never qualifies, §\ref{sec:2-3}.)

For known \textit{CGEL} members (\mention{lot, plenty, lots, bags, heaps, loads, oodles, stacks, remainder, rest, number, couple}₂), lexical identity is sufficient input for description: knowing a noun is on this list licenses comparison with the rest of the profile. The non-trivial projective test is candidate discovery, where the input criteria cannot simply repeat the targets.

For candidate new members, class discovery uses two diagnostics independently observable from any of the projected behaviours: quantificational semantics (an indefinite quantity rather than a precise count or kind) and high-frequency occurrence as head of an \mention{of}-complement construction. A candidate that satisfies both is then tested against the projected cluster:

\begin{itemize}
\item
  near-obligatory agreement override in the central cases, singular- or plural-direction depending on subgroup;\footnote{A pilot COCA check found every fully audited discriminating cell in the predicted direction. Counter-direction cells received more extensive manual filtering than predicted-direction cells, so the counts are attestational corroboration rather than estimates of grammatical rates. \mention{A lot of money} required parse-shift filtering, and the large \mention{a lot of people is/was} cell remains incompletely scrubbed. Full counts and methodology: Appendix~\ref{sec:A-3}.}
\item
  characteristic determiner patterns by subgroup (\mention{lot} most often with \mention{a}, also with \mention{the} in totality uses; \mention{plenty} typically without a determiner; the plural-only forms typically without a determiner; \mention{remainder, rest} typically with \mention{the}; \mention{number, couple}₂ typically with \mention{a});
\item
  partitive vs.~pseudo-partitive \mention{of}-complement availability;
\item
  behaviour under complement ellipsis: the \mention{of}-complement can be omitted while remaining understood, and it continues to determine the NP's number (\mention{A lot was spent on travel} vs \mention{A lot were arrested}; \textit{CGEL} [56], p.~349);\footnote{These bare uses are not \textit{CGEL}'s \term{fused-head} construction, where the head fuses with a determiner or modifier (\mention{Did you buy \textup{[}some\textup{]} yesterday?}; §9.1, p.~410). \textit{CGEL} analyzes bare QN NPs as ordinary heads with an elliptical complement (pp.~412, 503). The fusion concept returns, differently applied, in the fused relatives of §\ref{sec:2-5}.}
\item
  count/non-count percolation of the \mention{of}-complement to the whole NP (\mention{lot} \enquote{allows the number of the oblique to percolate up to determine the number of the whole NP}; \textit{CGEL}, p.~349), with consequences for selection by predicates that are count- or non-count-sensitive;
\item
  restricted premodification (\mention{a whole/huge lot} but not \mention{a remarkable plenty}).
\end{itemize}

The class's sub-group structure makes the cluster's predictions concrete. Table~\ref{tab:qn-matrix} displays the per-noun pattern for the central diagnostics. The entries are \textit{CGEL}-coded except in the two \mention{of}-complement columns, where the COCA pilot has revised them: the determined frame is attested for every plural-only member, and \mention{rest} and \mention{remainder} admit a bare complement when it denotes uniquely. The semantic classification of \mention{rest} remains open, and a fuller per-noun corpus audit is the within-domain check still needed for the profile claim.

{
\begin{longtable}[]{@{}
  >{\raggedright\arraybackslash}p{(\linewidth - 8\tabcolsep) * \real{0.20}}
  >{\raggedright\arraybackslash}p{(\linewidth - 8\tabcolsep) * \real{0.16}}
  >{\raggedright\arraybackslash}p{(\linewidth - 8\tabcolsep) * \real{0.19}}
  >{\raggedright\arraybackslash}p{(\linewidth - 8\tabcolsep) * \real{0.20}}
  >{\raggedright\arraybackslash}p{(\linewidth - 8\tabcolsep) * \real{0.25}}@{}}
\caption{Number-transparent quantificational nouns: behavioural matrix (\textit{CGEL}-based pilot).}\label{tab:qn-matrix}\\
\toprule\noalign{}
Sub-group / Noun & Determiner & Partitive \mention{of} N & Pseudo-partitive \mention{of} N & Override direction \\
\midrule\noalign{}
\endfirsthead
\toprule\noalign{}
Sub-group / Noun & Determiner & Partitive \mention{of} N & Pseudo-partitive \mention{of} N & Override direction \\
\midrule\noalign{}
\endhead
\bottomrule\noalign{}
\endlastfoot
\multicolumn{5}{@{}l}{\textbf{Singular, \mention{a}-determined}} \\
\mention{lot} & \mention{a}; also \mention{the} (totality) & $\checkmark$ (\mention{a lot of the money}) & $\checkmark$ (\mention{a lot of money}) & sg → pl with plural complement \\
\mention{number} & \mention{a} & $\checkmark$ (plural oblique) & $\checkmark$ (plural oblique) & sg → pl with plural complement \\
\mention{couple}₂ & \mention{a} & $\checkmark$ (plural oblique) & $\checkmark$ (plural oblique) & sg → pl with plural complement \\
\addlinespace
\multicolumn{5}{@{}l}{\textbf{Singular, usually without a determiner}} \\
\mention{plenty} & usually without a determiner & not typical & $\checkmark$ (\mention{plenty of money}) & sg → pl with plural complement \\
\addlinespace
\multicolumn{5}{@{}l}{\textbf{Plural, usually without a determiner}} \\
\mention{lots, bags, heaps, loads, oodles, stacks} & usually without a determiner & $\checkmark$, dispreferred; attested for every member (§\ref{sec:2-1}) & $\checkmark$ & pl → sg with non-count complement \\
\addlinespace
\multicolumn{5}{@{}l}{\textbf{Singular, \mention{the}-determined}} \\
\mention{remainder} & \mention{the} & $\checkmark$ (\mention{the remainder of the time}); bare whole only if bounded and unique (\ungram{\mention{the remainder of time}}) & n/a: not a quantity noun & sg/pl (complement-controlled) \\
\mention{rest} & \mention{the} & $\checkmark$ (\mention{the rest of the time}); bare whole only if bounded and unique (\mention{the rest of society}) & ? (classification disputed) & sg/pl (complement-controlled) \\
\addlinespace
\multicolumn{5}{@{}l}{\textbf{Boundary}} \\
\mention{majority, minority} & \mention{a} & $\checkmark$ & ? & variable \\
\mention{bunch} & \mention{a} & $\checkmark$ & $\checkmark$ & variable \\
\end{longtable}
}

One sentence per central sub-group shows the override directions the table codes.

\ea\label{ex:subgroup-sentences}
\ea \mention{A number of spots have appeared.}
\ex \mention{Plenty of people were stranded.}
\ex \mention{Heaps of money has been spent.}
\ex \mention{The rest of the money was returned.}
\z
\z

(\ref{ex:subgroup-sentences}a) and (\ref{ex:subgroup-sentences}c) repeat (\ref{ex:cgel-two-direction}): a singular \mention{a}-determined noun driven plural by its complement, a plural-only noun driven singular. (\ref{ex:subgroup-sentences}b) shows determinerless \mention{plenty} with the plural override; (\ref{ex:subgroup-sentences}d) shows the \mention{the}-determined sub-group, where direction is complement-controlled (plural instead with \mention{the rest of the people}). The directions in all four match the COCA pilot's counts in Table~\ref{tab:override-audit}.

\textit{CGEL} treats override as obligatory in the central cases; the pilot corroborates the direction in selected cells but does not audit every class member or estimate rates. Complement-ellipsis behaviour and count/non-count percolation remain \textit{CGEL}-based predictions rather than per-noun corpus results. Premodification is restricted across the class, though the restriction looks like one on modifier type rather than on premodification as such: degree and size modifiers are available (\mention{a whole/huge lot}; \mention{huge}, \mention{enormous}, and \mention{vast piles of money}), while evaluative ones are marginal (\mention{a remarkable plenty}), and \mention{plenty} takes minimal modification generally.

Each sub-group covaries: knowing a candidate noun's determiner pattern and partitive availability narrows the sub-group, which then licenses the override-direction prediction (and vice versa). One qualification on the partitive diagnostic. For \mention{remainder}, the absence of a quantity reading follows from what kind of noun it is; for \mention{rest}, the lexical-semantic classification remains unresolved. The diagnostic therefore carries less weight than the determiner and override behaviours (Appendix~\ref{a.2-partitive-vs.-pseudo-partitive-of-complements}). The boundary cases have weaker covariance, which is why §\ref{sec:2-2}'s gradient framing applies to them.

Table~\ref{tab:override-audit} displays the raw agreement-direction counts. The discriminating cells are the informative ones, where the head's number and the complement's diverge. The matched cells confirm the expected direction but cannot distinguish the override from head-driven agreement. Counter-direction cells received more extensive manual checking than predicted-direction cells, so the table reports an audit trail rather than estimated rates. The large \mention{a lot of people} cell remains asymmetric, and the smallest cells remain feasibility checks. Referential \mention{the number of} is a negative control: the same noun agrees head-driven under \mention{the} and overrides only under quantificational \mention{a} (Appendix~\ref{sec:A-5}), so the contrast belongs to the construction, not the lexeme.

\begin{table}
\centering
\small
\caption{Raw agreement-direction counts in the COCA pilot. Counts are after the filtering described in Appendix~\ref{sec:A-3}; the audit-status column records the evidential limitation.}\label{tab:override-audit}
\begin{tabular}{@{}p{0.22\linewidth}p{0.12\linewidth}rrp{0.38\linewidth}@{}}
\toprule
Subject & Predicted & Plural & Singular & Audit status \\
\midrule
\mention{a lot of people} & plural & 4,195 & 85 & Large counter-direction cell not symmetrically scrubbed \\
\mention{a lot of money} & singular & 1 & 90 & Counter-direction cell checked; parse-shift filtering \\
\mention{a number of people} & plural & 100 & 0 & False positives checked \\
\mention{lots of people} & plural & 348 & 0 & False positives/non-standard uses checked \\
\mention{lots of money} & singular & 0 & 24 & Small cell; false positive checked \\
\mention{plenty of people} & plural & 79 & 0 & Small cell \\
\mention{plenty of money} & singular & 0 & 6 & Small cell \\
\mention{the rest of the people} & plural & 19 & 0 & Small cell \\
\mention{the rest of the money} & singular & 0 & 19 & Small cell \\
\bottomrule
\end{tabular}
\end{table}

This is the prospective design for a new-member case: quantificational semantics and \mention{of}-complement-head distribution are inputs, while the cluster of agreement, determiner, partitive, ellipsis, count/non-count, and premodification behaviours supplies the targets. Boundary cases like \mention{majority} and \mention{minority} (§\ref{sec:2-2}) and the dual-sense \mention{couple}₁ (collective) versus \mention{couple}₂ (number-transparent) show graded membership. The design distinguishes prediction from redescription, but the present evidence supports only selected targets.

The more demanding test asks what the profile predicts beyond the in-grammar cluster. Diachronic entry and cross-linguistic analogues are plausible targets, but neither is tested here. One candidate-member test has been run. \mention{Piles} was admitted on two criteria only, quantity meaning and frequent \mention{of}-complement distribution, neither of which is among the predicted behaviours. The determiner-frame prediction receives partial support: \mention{a piles of} returns no COCA hits, as the plural-only subgroup requires, while the determined frame \mention{the piles of} N selects the literal reading in the available sample. This result separates the two senses of \mention{piles}, but it does not validate the full profile.

A follow-up sweep suggests that the separation extends across the plural-only subgroup. The determined frame recruits the literal homonym wherever a member has one: \mention{stacks of the *} is dominated by library stacks and smokestacks, \mention{heaps of the *} by literal and colliery heaps, and \mention{bags} and \mention{loads} by containers and cargo. \mention{Oodles}, the one member with no literal homonym, returns quantificational hits exclusively. This is a useful subgroup hypothesis, not a completed projection test.

The agreement prediction couldn't be tested at all. COCA has 91 tokens of \mention{piles of money} and one followed by a finite form of \mention{be}; adding COHA's two centuries brings the discriminating non-count cell to three tokens, one singular and two plural. Nothing there supports the override and three tokens cannot weigh against it either. A container-noun control (\mention{sacks/boxes/crates of money}) returns nothing at all, and the left contexts of the 91 tokens locate the problem precisely: 41 have \mention{piles of money} as the object of a verb (\mention{make, spend, earn, amass, inherit}) and 13 as the complement of a preposition, so at least 59\% sit in complement position, while the single token preceded by finite \mention{be} is copular or existential, with agreement running from a different subject. Quantity expressions of this shape occur as complements, not as subjects of finite \mention{be}. The same sparsity blocked the inanimate \mention{bunch} contrast in §\ref{sec:2-2}.

So the override prediction, which is the profile's headline, is corpus-testable only for members that appear in subject position, and candidate-member discovery for this class will need acceptability judgments rather than frequency. The \mention{piles} result is therefore partial prospective support plus a feasibility failure. The protocol, full returns, and corrections the wildcard queries forced on Appendix~\ref{a.2-partitive-vs.-pseudo-partitive-of-complements} are in the supplementary material (see Data availability).

The cluster makes number-transparent quantificational nouns the strongest candidate \term{projectibility profile} examined here: an observable base supports projection to a structured set of targets, with conventional lexical inheritance as the proposed stabilizer. The profile is a candidate for Slater's \autocite*{slater-2015-natural-kindness} stability-plus-projection diagnostic, but the within-domain check beyond selected agreement and determiner behaviour remains open. A profile prediction would fail if a noun satisfying the input diagnostics systematically lacked the target cluster. The status of \mention{rest} is an explicit analytical fork: it remains an access-transparency case under either treatment, but if it is a part/subset expression rather than a quantity noun, it should not count as input to the quantity-QN candidate-discovery generalization. Whether the profile earns a local-kind label is a downstream question for §\ref{sec:7}.

\subsection{Transparent free relatives}\label{sec:2-5}

A second morphosyntactic-transparency case fits the schema. Transparent free relatives (TFRs) are a sub-type of fused-relative construction (\textit{CGEL}, ch.~12, §6, pp.~1068ff.) in which the matrix accesses properties of an embedded predicate (the \term{nucleus}) through \mention{what}: \mention{what they call ecstatic}, \mention{what we regard as essential}, \mention{what seems to be a tourist}. \term{Fused} is \textit{CGEL}'s term for one word jointly realizing two functions: in a fused relative, \enquote{the head of an NP fuses with the relativised element in a relative clause} (p.~1073), so \mention{what} in \mention{what she wrote} is at once the NP's head and the relativized element. This paper uses TFRs only as a second profile case. The analysis followed here is Kim's construction-grammar one, developed across \textcite{kim-2011-tfrs}, \textcite{kim-2017-free-relatives}, \textcite[ch.~6.8]{kim-2026-form-function}, and its revision in \textcite[§5]{kim-2026-indeterminacy-free-relatives}, with the corpus material in \textcite{kim-reynolds-2026-tfr-two-kinds}. The verdict below is conditional on that analysis: rival treatments locate the mediator elsewhere, and where the mediator sits is exactly what the schema needs fixed, so a reader who prefers Wilder's \autocite*{wilder-1999-transparent-free-relatives} parenthetical-with-backward-deletion account, van Riemsdijk's \autocite*{vanriemsdijk-2006-free-relatives} graft analysis, or Grosu's \autocite*{grosu-2003-unified-theory-tfrs} unified treatment should expect the slots to be filled differently: each locates the mediator, and so the access path, somewhere else. The verdict is therefore analysis-relative: positive under Kim's constructional analysis, indeterminate on an analysis-neutral reading, since the competing treatments aren't adjudicated here.

\ea\label{ex:tfr-r-relations}
\ea \mention{What appear to be trousers are on the chair.}
\ex \mention{What they call essential workers are exhausted.}
\ex \mention{He was utterly penniless and what they call dimwitted.}
\z
\z

TFRs provide a second, conditional grammatical illustration, but the evidence is deliberately bounded. They support two relations at different levels of the construction. One is agreement: matrix agreement accesses the nucleus's number profile through singular \mention{what}. The other is matrix selection, whose accessible property is the nucleus's syntactic category; \term{category flexibility} is the outcome, that a non-NP nucleus can satisfy the matrix's requirement, as in (\ref{ex:tfr-r-relations}c), where the nucleus is an AP. Predicative position admits both NPs and APs, so that example shows the nucleus's category rather than establishing a category-discriminating environment. Kim \autocite*[36--37]{kim-2026-indeterminacy-free-relatives} treats the second as the marked sub-case of a wider \term{selection transparency}, on which the matrix's selectional requirement is satisfied by the nucleus generally, NP nuclei included, as when a preposition selects the nucleus through the \mention{what}-clause (\mention{he speaks in what linguists call a Northern dialect}).

The concealing counterpart for parent-level agreement is the ordinary fused relative, where \mention{what} itself controls agreement. That control is family-level rather than minimally matched: an ordinary fused relative like \mention{what she wrote} lacks a nucleus, so it can't hold the concealed property constant while varying only the access. It establishes that the fused-relative family admits mediator-controlled agreement, not that a matched variant conceals this nucleus's number. The fourth gate requires the counterpart's evidential strength to be stated, and it enters here at that weaker, family-level grade.

The pilot reported with the joint analysis supplies a bounded check on the category contrast. In a full coding of 485 archived \mention{call} hits, 185 were genuine AP-nucleus TFRs, 144 of them externally predicative, with one token coordinated into attributive position. Exhaustive checks of the candidate sets for right-edge non-NP nuclei with \mention{seem} or \mention{appear} (721 tokens across eight queries, run or re-verified against COCA in August 2026) found no adjectival or adverbial case and a single PP-nucleus token in supplement position; the working paper reports the query details. This is distributional evidence, not a completed acceptability study; the relevant category-discriminating environments and judgments remain to be added to the joint paper.

The observable base is constructional sub-type plus verb-class membership, and the examples instantiate both classes Kim \autocite*[40]{kim-2026-indeterminacy-free-relatives} distinguishes: the \term{evidential} type (\mention{seem~to~be, appear~to~be, happen~to~be}; \ref{ex:tfr-r-relations}a), which takes NP nuclei only in the current analysis, and the \term{attributional} type (\mention{call, consider, describe~as, regard~as, take~to~be}; \ref{ex:tfr-r-relations}b,~c), which licenses AP, AdvP, PP, and \mention{-ing} nuclei as well. On the inheritance hierarchy Kim \autocite*[46]{kim-2026-indeterminacy-free-relatives} sets out, the parent licenses agreement, selection, binding, and coordination transparency across both classes, and the attributional daughter adds category flexibility. \enquote{Coordination transparency} there is Kim's name for a TFR diagnostic, that the construction conjoins as its nucleus's category; it doesn't reverse §\ref{sec:2-3}'s verdict on ordinary coordinate NPs, which concerns a different configuration. The pilot therefore supports a provisional parent/daughter contrast, with constructional inheritance as the proposed stabilizer. The full verb-class inventory, AP/PP stratification, and rival analyses remain part of the joint paper.

\subsection{Bridge}

The cases in Section \ref{sec:2} share a relation: morphosyntactic matching mediated by a construction or lexical class. Section \ref{sec:3} marks form--meaning applications at the boundary; the cross-linguistic reach of the schema, including the controlled Spanish and Tsez cases and the topology survey, is Section \ref{sec:6}'s topic.

\section{Boundary applications beyond grammar}\label{sec:3}

The schema can also describe form--meaning access when $R$ is interpretation, learning, or processing rather than agreement. Constructional compositionality \autocite{goldberg-2006-constructions-at-work}, rated transparency in noun--noun compounds \autocite{bell-schafer-2016-modelling-transparency}, and masked morphological priming \autocite{heyer-kornishova-2018-priming-eventually} supply possible applications. In each case, the analyst would need to identify a bearer, a distinct mediator, a relation, and a constructional counterpart that could conceal the relevant property.

These examples mark scope rather than deliver additional worked profiles. Their relations, items, diagnostics, and proposed stabilizers differ, and the available studies don't establish a common within-items cluster. Constructional interpretation, compound ratings, and morphological priming belong to a research programme outside the evidential spine of this paper. The examples show why the schema keeps $R$ explicit: a construction may be accessible for one relation and opaque for another. They don't warrant a cross-domain projectibility claim here.

\section{Typological transparency}\label{sec:6}

This section is methodological. The previous sections worked through cases in English; §\ref{sec:3} marks form--meaning applications without using them as evidence for a cross-domain profile. The question here is how to move from language-particular cases to a cross-linguistic claim. Haspelmath's \autocite*{haspelmath-2010-comparative-concepts} descriptive-versus-comparative-concept distinction is the right discipline.

\subsection{Haspelmath's distinction}

Haspelmath \autocite*{haspelmath-2010-comparative-concepts} argues that crosslinguistic comparison shouldn't assume universally available descriptive categories\footnote{Throughout, \mention{category} is used in Haspelmath's \mention{descriptive category} sense: a language-particular grammatical grouping fixed by a language's own distributional criteria, and neutral on whether that grouping is projectible. This differs from the sense used elsewhere in this research programme, on which a \mention{category} just is a projectible kind \autocite{reynolds2025hpcbook}; that kind-level notion is carried here by the \mention{projectibility profile} (§\ref{sec:7-1}), not by \mention{category}.} like ADJECTIVE, PASSIVE, ACCUSATIVE, or SUBJECT. Each language has its own descriptive categories, fixed by the language-particular criteria that make them work in that language \autocite[664]{haspelmath-2010-comparative-concepts}. Comparative concepts, by contrast, are analyst-created tools for comparison; they aren't part of any language system and can only be more or less well suited to the comparative task, rather than right or wrong \autocite[665]{haspelmath-2010-comparative-concepts}. They have to be built from universally applicable building blocks: conceptual-semantic notions, general formal notions, and other comparative concepts.

One pair makes the relation concrete. Haspelmath's comparative concept \enquote{dative case} is defined semantically, as a marker among whose functions is coding the recipient of a physical-transfer verb \autocite[666]{haspelmath-2010-comparative-concepts}. Language-particular categories then match it or not independently of their traditional names: the Finnish Allative matches, while the Nivkh Dative-Accusative doesn't, despite the name, because it codes a causee rather than a recipient \autocite[665--666]{haspelmath-2010-comparative-concepts}. Categorial particularism, Haspelmath's preferred position, holds that descriptive categories are language-particular and that comparison runs through analyst-defined comparative concepts standing in many-to-many relations with them. The position it rejects, categorial universalism, assumes a substantial set of crosslinguistic descriptive categories from which languages select \autocite[663]{haspelmath-2010-comparative-concepts}; \textcite{newmeyer-2007-crosslinguistic-formal-categories} is its canonical statement, taken up in §\ref{sec:6-objection}.

The methodological consequence is straightforward. A typological generalization isn't a claim about a descriptive category instantiated in every language. It's a claim about how analyst-defined comparative concepts are realized in language-particular ways. That nominative case-marking is widespread doesn't entail that \enquote{nominative} picks out the same descriptive category in any two languages.

\subsection{The umbrella schema as a comparative concept}

The umbrella schema is a comparative concept, not a descriptive category. It's built from the materials Haspelmath \autocite*[665]{haspelmath-2010-comparative-concepts} admits: \mention{property} and \mention{relation} are conceptual-semantic primitives, \mention{form} is a general formal notion. Asking the schema's question of a language isn't claiming that the language has a category corresponding to \term{transparency}.

\mention{Accessible through} is the contested slot, since accessibility has discipline-specific senses (discourse-pragmatic salience, psycholinguistic retrievability, formal-semantic substitution-licensing). The schema's sense is more abstract: $P$ is accessible to $R$ through $X$ when $R$ tracks $P$ across $X$ rather than tracking a property of $X$ itself. The case studies operationalize that differently, by agreement targeting, constructional interpretability, rated transparency, or priming magnitude, but the comparative concept lives above the operationalizations.

The non-instance test (§\ref{sec:2-3}) is what keeps the concept honest, because its counterpart is fixed language-internally. The test asks whether the language under study attests a same-type variant of $X$ whose structure makes $P$ inaccessible to $R$. For English number-transparent NPs that counterpart is head-driven agreement, and the transparent class is the marked exception against it; the pattern isn't exported. The procedure is identical in any language: identify $X$, $P$, and $R$ by that language's own distributional criteria, then look for the concealing variant. Where it exists, the transparent case is informative; where the construction-type admits none, as with English coordination, the schema returns a non-instance. So \mention{accessible through} imports no English-particular grammar. It inherits whatever counterpart the local construction-type supplies.

So the bindings are language-particular while the question isn't. In English, $X$ is the quantificational noun plus \mention{of}-complement construction, a \textit{CGEL} descriptive category; $P$ is the complement's number/count profile; $R$ is subject-verb agreement. In the worked cases below, Spanish partitive subjects and Tsez complement clauses fill the slots by different language-particular criteria. That's what makes the typological question tractable: not \enquote{do all languages have number-transparent constructions?}, which would look for English-style NPs everywhere, but \enquote{which configurations in this language let a property remain accessible to a relation that the surface form might have hidden?} The answers are expected to differ, and the English, Spanish, and Tsez cases aren't instances of one descriptive category.

Two limits on what that buys. The claim at the comparative-concept level is structural, that the concept lives above the descriptive category whatever the language, rather than a demonstrated comparative result: the extensions on hand are the Spanish and Tsez cases below. The form--meaning family (§\ref{sec:3}) isn't itself a comparative concept; the schema is, and the family is a set of language-particular phenomena and measurements compared through it.

\subsection{Cross-linguistic cases and topologies}\label{sec:6-topologies}

Two worked languages let the question be asked outside English, one of them outside the agreement type English uses. Spanish has an analogous agreement-access pattern. In a partitive subject like \mention{la mayoría de los manifestantes} (`the majority of the demonstrators'), the finite verb may agree in the plural with the complement or in the singular with the quantificational noun \mention{mayoría}, as in (\ref{ex:spanish-mayoria}).

\ea\label{ex:spanish-mayoria}
\ea \gll La mayoría de los manifestantes gritaban. \\
the majority of the demonstrators shout.\textsc{ipfv.3pl} \\
\glt `The majority of the demonstrators were shouting.'
\ex \gll La mayoría de los manifestantes gritaba. \\
the majority of the demonstrators shout.\textsc{ipfv.3sg} \\
\glt `The majority of the demonstrators was shouting.'
\z
\z

The plural (\ref{ex:spanish-mayoria}a) is the more frequent option and the only one available under a plural predicative \autocite{rae-concordancia-mayoria}. Both variants are attested on the one subject, much as English \mention{majority} and \mention{minority}, \mention{mayoría}'s boundary cognates, take occasional singular agreement alongside the dominant plural override (§\ref{sec:2-2}). Spanish has its own literature on the alternation, which treats it as \term{concordancia ad sensum} and reaches two conclusions that matter here. First, both options are genuine agreement rather than agreement displaced by government, so the plural isn't a species of \term{silepsis} \autocite{sanjulian2018-cuantificadores}. Second, the alternation isn't free across syntactic contexts, and singular agreement with the quantifier noun is compatible with a distributive reading of the clause \autocite{sanjulian2018-cuantificadores}, which blocks the inference that notional plurality forces the override. The partitive and pseudo-partitive structures the alternation runs on are themselves heterogeneous in Spanish \autocite{sanjulian2018-pseudopartitivas}, as they are in English (§\ref{sec:2-1}).

What travels to Spanish is the access pattern, not the English cluster: no Spanish determiner, complementation, ellipsis, or premodification profile has been shown here. And the concealment control is weaker than the English one, since the singular option with \mention{mayoría} is dispreferred rather than blocked.

Tsez (Nakh-Daghestanian; head-final and ergative; about seven thousand speakers) gives a case in a different marking type. Tsez verbs prefix-agree in noun class with their absolutive argument, so the agreement is marked on the verbal head. A clausal argument is class IV, the class that includes clauses, so a matrix verb's default is class IV agreement with the complement clause as a unit; but when the embedded absolutive is a topic, the matrix verb can instead agree with it \autocite[436]{potsdam-polinsky-1999}.

\ea\label{ex:tsez-lda}
\ea \gll enir [už-ā magalu b-āc'ru-łi] r-iyxo \\
mother [boy-\textsc{erg} bread.\textsc{iii.abs} \textsc{iii}-ate-\textsc{nmlz}] \textsc{iv}-knows \\
\glt `The mother knows the boy ate the bread.' (default: the matrix verb is class IV, agreeing with the clause)
\ex \gll enir [už-ā magalu b-āc'ru-łi] b-iyxo \\
mother [boy-\textsc{erg} bread.\textsc{iii.abs} \textsc{iii}-ate-\textsc{nmlz}] \textsc{iii}-knows \\
\glt `The mother knows the boy ate the bread.' (long-distance agreement: the matrix verb is class III, agreeing with embedded \mention{magalu} `bread')
\z
\z

In the schema's terms, $X$ is the embedded clause, $P$ is the noun class of the embedded absolutive, and $R$ is matrix verb agreement. The concealing variant, and here genuinely the default, is the class IV agreement in (\ref{ex:tsez-lda}a), attested in the same language and the same sentence, and (\ref{ex:tsez-lda}b) is the transparent override. The counterpart isn't \enquote{the nearest head's own feature}, as with English number-transparent nouns, but \enquote{default class IV for a clausal argument}: a different default in a different language type. That's what the comparative concept predicts, the question constant and the language-particular answer not. Whether the override is analysed as covert movement or as long-distance agreement is a framework-internal question the schema brackets, as with extraction in §\ref{sec:2}; the fuller analysis is in \textcite{polinsky-potsdam-2001}.

Spanish and Tsez are both outward endpoint cases: a property associated with material inside $X$ reaches an external agreement relation. The same schema maps several other accessibility topologies, provided each remains subject to the §\ref{sec:2-3} counterpart test: identify $X$, $P$, and $R$, then ask what same-language pattern would make the proposed access disappear. Appendix~\ref{sec:B} sets out four of them with their applicability and status conditions (null-source subpart access in Aleut, inward dependent access in Lardil and Kayardild, path-distributed access in Chamorro, and the feature-relay boundary case in Kutenai and Algonquian obviation), along with a targeted check of head-marked systems.

Two results from that survey bear on the argument here. Aleut is a qualified positive: the controller of $R$ can be recoverable as a subpart of a non-subject NP without being overt, though Merchant's analysis treats the unpronounced possessor as syntactically present pro, which keeps the case distinct from overt intact-NP transparency \autocite[195]{merchant-2011-aleut-case-matters}. And the head-marked candidates cluster around possessors, most of them analysed as externalized, by possessor raising, applicative mediation, or external possession, rather than as agreement with a dependent that stays NP-internal \autocite{ritchie-2016-prominent-internal-possessors,deal-2017-external-possession}. So the diagnostic point is narrow: overt sensitivity to a possessor isn't enough, and the decisive question is whether an embedded dependent remains NP-internal while controlling clause-level morphology.


\subsection{Mapping transparency and access transparency}\label{sec:6-mapping}

There's an established typological programme under the name \term{transparency}, and it classifies the paper's paradigm cases as non-transparent. Setting the two side by side is the clearest available demonstration that \term{transparency} isn't one projectible kind.

Leufkens defines transparency as \enquote{a one-to-one relation between meaning and form}, treating every structure that violates such a relation as non-transparent \autocite[2]{leufkens2015}, and ranks 22 languages by how many non-transparent features their grammars show; the framework is extended to a larger sample and an implicational hierarchy in \textcite{hengeveld-leufkens-2018}. Opacity divides into four categories according to how the one-to-one relation fails \autocite[16, §§4.1--4.4]{leufkens2015}: redundancy, discontinuity, fusion, and form-based form. Agreement is the leading instance of the first. Where one unit of meaning is expressed by several forms, \enquote{one of the forms is not necessary; it is redundant, as it does not provide any additional information}, and clausal and phrasal agreement are catalogued in that category \autocite[18, §§4.1.3--4.1.4]{leufkens2015}.

That's the opposite verdict from this paper's on the same configurations. Subject--verb agreement is the paradigm $R$ for which a property remains accessible through a mediator; on the one-to-one conception it's redundant marking, and a language with more of it is less transparent. Both verdicts are correct within their own question. They're different questions: whether a language's levels map one-to-one, and whether a particular relation can reach a property across a mediator that would otherwise conceal it. Call the first \term{mapping transparency} and the second \term{access transparency}.

English number-transparent QNs cross-classify. They're an access-transparency instance. Being agreement, they're also a mapping-transparency violation. A language could reduce its agreement, becoming more transparent in Leufkens's sense, and thereby lose the very relations that make access transparency visible. The divergence isn't a disagreement to be resolved by defining terms more carefully, because the two programmes individuate by different relata: levels of organization in one, mediator-property-relation triples in the other. An umbrella that delivered opposite verdicts on one construction while both verdicts stood wouldn't be a kind, and that's what §\ref{sec:7} concludes on independent grounds.

Two further payoffs. Leufkens's inventory is organized, so it supplies a principled comparison class rather than another loose sense of the word, and the redundancy/discontinuity/fusion/form-based-form taxonomy is a checklist this paper's cases can be run against. And her finding that every language sampled has some non-transparent features, most often of the redundancy type \autocite{leufkens2015}, means agreement relations are common, so they're where a search for mediated access is most likely to find candidates. That says nothing about how often the search succeeds.

\subsection{An objection and a reply}\label{sec:6-objection}

A reader might object that the umbrella schema looks like a universal descriptive category smuggled in as a comparative concept. The sharpest form of the objection targets the non-instance test: \mention{accessible through} presupposes a concealing counterpart, and for the English cases that counterpart is head-driven agreement, which is independently the unmarked pattern, a language-particular fact. If the schema's deepest primitive is parasitic on an English-shaped default, then calling it universally applicable looks like categorial universalism in disguise.

The reply isn't that the schema is \enquote{just a question}: a universal category can be framed as a question too, so that move wouldn't tell a comparative concept apart from a category in disguise. The reply is that the counterpart the objection points to is fixed \mention{language-internally}, which is how comparative concepts are meant to work. The concealing counterpart is not assumed to be English head-driven agreement; it's whatever same-type variant the language under study attests, by the procedure set out earlier in this section. So the schema doesn't import an English default; it inherits each language's own.

Three things follow that a universal descriptive category wouldn't deliver. The variables are bound by language-particular distributional criteria, not by a fixed cross-linguistic definition. The schema stands in many-to-many relations with descriptive categories: one language's category can fill $X$, $P$, or $R$ under different bindings, and several categories can realize one binding. It also carries explicit status conditions: a candidate is a non-instance where the construction-type lacks the required mediator or bearer, and a profile claim is demoted where the predicted access fails to project, as in the \mention{bunch} localization of §\ref{sec:2-2}. The honest residue is that the test is only as sharp as the prior individuation of construction-type. Deciding what counts as the same type with and without concealment is itself a language-internal analysis, so the diagnostic is crisp where that individuation is independently motivated and softer where it isn't. A candidate counterpart that isn't independently motivated can't supply a confident positive. The procedure's refusal to overgeneralize is a useful stress test, though indeterminacy alone doesn't validate the test. Haspelmath's \autocite*[665]{haspelmath-2010-comparative-concepts} criterion for a good comparative concept, that it can be more or less well suited to the task of permitting crosslinguistic comparison, is exactly the standard the umbrella schema is meant to meet.

Newmeyer's specific argument deserves a direct response. His case is that typological generalizations must reference not just a universal concept but the specific form realizing it: negative particles and negative auxiliaries express the same concept, but only the latter pattern with verbs in word-order typology, so a generalization stated over \enquote{negation} misses the regularity \autocite[136]{newmeyer-2007-crosslinguistic-formal-categories}. The point is correct on its own terms. What follows from it is weaker than crosslinguistic \mention{descriptive} categories, because the formal reference it demands enters at the $X$ slot, and the slot and its values are different things. The slot is \mention{typed} with general-formal notions in Haspelmath's \autocite*[666]{haspelmath-2010-comparative-concepts} sense: construction, clause, NP, dependency domain. A value of the slot is a language-particular configuration fixed by that language's criteria, the number-transparent QN construction in English and the configurations the analyst identifies in each comparison language. The same holds for $P$ and $R$: the slots are abstractly typed, while their values are particular number categories, grammatical relations, and agreement systems. Nothing here requires equating configurations token-for-token across languages. The concealing counterpart that gives the non-instance test its grip is itself such a general-formal concept, not an imported universal category.

The metaphysical conclusion in §\ref{sec:7}, that the umbrella schema is a tool for locating projectibility profiles rather than itself a kind, is consistent with this. Three levels stay distinct: the comparative concept (mediated accessibility) lives at the cross-linguistic methodological level; descriptive categories (English number-transparent quantificational nouns, Spanish partitive subjects, Tsez long-distance agreement, and so on) live at the language-particular grammatical level; \term{projectibility profiles} are observable bases licensing non-trivial projection, stabilized by identifiable sources, realized within or across such descriptive categories. Kindhood is a downstream verdict on profiles that earn it. The paper doesn't collapse the levels.

Two limits on the position taken here. Haspelmath's distinction is itself contested and this paper doesn't adjudicate; it's adopted as the safer default, since it requires no commitment to descriptive-category universality, and a reader more comfortable with categorial universalism can take the schema as a working concept whose status is secondary to its methodological pay-off. And Croft's \autocite*{croft-2001-radical-construction-grammar} Radical Construction Grammar would route the same commitment through construction-similarity hierarchies rather than comparative concepts; the schema is meant to sit above any single construction-grammar variant, so that route is left open rather than argued against.

One consequence carries into §\ref{sec:7}: cross-linguistic projection at the profile level isn't entailed by cross-linguistic projection at the comparative-concept level. Asking the question everywhere doesn't make the answers travel.

\section{Projectibility profiles}
\label{sec:7}

Section 1 made projectibility a secondary question. A positive access relation may support induction, but the comparative concept itself carries no predictions. This section records which grammatical profiles support further projection and what appears to stabilize them. The primary field is syntax, with targets such as agreement, syntactic-category licensing, and complement selection.

The distinction matters because \term{access transparency} and \term{projectibility profile} answer different questions. Applicability asks whether a case fills the schema's roles and controls. A relational verdict asks whether the relation reaches the bearer across the mediator. A profile verdict asks whether an observable base supports projection to further behaviour. The three verdicts should travel separately.

\subsection{Applicability and profile verdicts}\label{sec:7-1}

The umbrella schema is a comparative concept, not a single profile. Its grammatical applications vary in $X$, $P$, and $R$: QNs are the primary demonstration and TFRs a conditional grammatical illustration; Spanish and Tsez are controlled extensions; the Appendix~\ref{sec:B} cases are candidate topologies or boundary cases. No common stabilizer licenses projection across all of them. The schema's contribution is the four-gate applicability protocol and its three-way verdict, not a prediction about every case called \mention{transparent}.

Table 3 assigns the roles and controls case by case. A blank counterpart marks a missing control rather than a negative empirical result. The final column uses the terminology of the protocol: positive, non-instance, or indeterminate.

Table 3. Roles and controls filled, case by case.

{\def\LTcaptype{none}\footnotesize
\begin{longtable}[]{@{}
  >{\raggedright\arraybackslash}p{(\linewidth - 8\tabcolsep) * \real{0.17}}
  >{\raggedright\arraybackslash}p{(\linewidth - 8\tabcolsep) * \real{0.21}}
  >{\raggedright\arraybackslash}p{(\linewidth - 8\tabcolsep) * \real{0.17}}
  >{\raggedright\arraybackslash}p{(\linewidth - 8\tabcolsep) * \real{0.23}}
  >{\raggedright\arraybackslash}p{(\linewidth - 8\tabcolsep) * \real{0.22}}@{}}
\toprule\noalign{}
Configuration & Property, and its bearer & Mediator & Relation (diagnostic) & Concealing counterpart \\
\midrule\noalign{}
\endhead
\bottomrule\noalign{}
\endlastfoot
QN + \mention{of}-complement & number/count profile, borne by the complement & the QN, morphologically mismatched & agreement (verb form) & head-driven agreement \\
\addlinespace
TFR, parent & number profile, borne by the nucleus & \mention{what} & agreement (verb form) & ordinary fused relative \\
\addlinespace
TFR, attributional daughter & syntactic category, borne by the nucleus & \mention{what} & matrix selection (category flexibility) & evidential daughter, NP nuclei only \\
\addlinespace
Collective NP & member plurality, borne by the head itself & -- (same sign) & agreement (verb form) & -- \\
\addlinespace
Coordination & summed number, borne by no single conjunct & -- & agreement (verb form) & -- \\
\addlinespace
Deferred reference & number, borne by the designatum & the subject NP & agreement (verb form) & -- \\
\addlinespace
Spanish partitive subject & number, borne by the complement & \mention{mayoría} & agreement (verb form) & singular option, dispreferred \\
\addlinespace
Tsez complement clause & noun class, borne by the embedded absolutive & the clause & matrix agreement (class prefix) & class IV, the clausal default \\
\addlinespace
Aleut (App.~\ref{sec:B}) & person/number, borne by an unpronounced possessor & the possessed NP & anaphoric verb morphology plus subject case & overt-possessor alternant \\
\addlinespace
Lardil/Kayardild (App.~\ref{sec:B}) & clausal TAM, borne by the clause & the dependency domain & dependent-nominal morphology & -- \\
\addlinespace
Chamorro (App.~\ref{sec:B}) & grammatical relation of the extractee & the extraction path & path morphology & -- \\
\addlinespace
Obviation (App.~\ref{sec:B}) & possessor obviation, borne by the possessor & the possessed NP & predicate tracks the head's derived feature & -- \\
\end{longtable}
}

Table 4. Applicability, relational, and profile verdicts.

{\def\LTcaptype{none}
\begin{longtable}[]{@{}
  >{\raggedright\arraybackslash}p{(\linewidth - 4\tabcolsep) * \real{0.25}}
  >{\raggedright\arraybackslash}p{(\linewidth - 4\tabcolsep) * \real{0.39}}
  >{\raggedright\arraybackslash}p{(\linewidth - 4\tabcolsep) * \real{0.36}}@{}}
\toprule\noalign{}
Case & Applicability and relational verdict & Profile verdict \\
\midrule\noalign{}
\endhead
\bottomrule\noalign{}
\endlastfoot
Number-transparent QNs (§\ref{sec:2-1}) & Positive: complement number/count reaches agreement through the QN & Strongest candidate; selected agreement and determiner tests support projection, but the full cluster remains open \\
\addlinespace
Transparent free relatives (§\ref{sec:2-5}) & Positive conditional on Kim's constructional analysis: nucleus properties reach agreement and selection through \mention{what} & Second syntactic candidate; pilot verb-class evidence, not a completed profile audit \\
\addlinespace
Collective NPs (§\ref{sec:2-2}) & Non-instance: bearer and mediator are layers of one sign, not distinct elements & None claimed \\
\addlinespace
Ordinary coordination (§\ref{sec:2-3}) & Non-instance: no mediator conceals the summed features & None \\
\addlinespace
Deferred reference (§\ref{sec:2-3}) & Non-instance: no bearer occurs within $X$ & None \\
\addlinespace
Spanish partitive subjects (§\ref{sec:6-topologies}) & Positive with a weak counterpart: complement number reaches agreement despite a dispreferred singular option & No English-like profile shown \\
\addlinespace
Tsez long-distance agreement (§\ref{sec:6-topologies}) & Positive: embedded noun class reaches matrix agreement against a class-IV counterpart & None claimed \\
\addlinespace
Aleut (§\ref{sec:B}) & Positive with qualification under Merchant's analysis of a null possessor & None claimed \\
\addlinespace
Lardil/Kayardild (§\ref{sec:B}) & Indeterminate: the inward relation is described, but a same-family counterpart hasn't been established & None claimed \\
\addlinespace
Chamorro (§\ref{sec:B}) & Indeterminate: path morphology is a candidate relation, without the required counterpart & None claimed \\
\addlinespace
Obviation (§\ref{sec:B}) & Non-instance or boundary case: the predicate tracks a feature relayed onto the head & None claimed \\
\end{longtable}
}

The table separates a missing control from a failed relation. A case can be positive while its profile remains untested, as with Tsez and the Spanish partitive. A proposed profile can also weaken while the local grammatical relation survives, as the sparse inanimate \mention{bunch} result shows for the QN discussion. Indeterminate cases shouldn't be promoted to positive instances merely because their topology is suggestive.

\subsection{What the primary profiles project}\label{sec:7-2}

The QN class is the strongest candidate because its observable base is independently available. Known members form a small conventional inventory, while candidate members can be admitted on quantificational semantics and frequent \mention{of}-complement distribution before their projected behaviours are examined. Determiner frames, complement selection, and agreement direction covary by subgroup. The COCA audit supports selected directions, but asymmetric filtering and sparse subject-position data leave the end-to-end cluster open.

The \mention{piles} test illustrates the proper evidential grade. Its determiner prediction received partial prospective support and revealed a broader literal-versus-quantificational sense split across the plural-only subgroup. Its agreement prediction wasn't testable because subject-position tokens were too sparse. The result supports candidate status and identifies the next feasible method, probably acceptability judgments, rather than validating the whole profile.

The status of \mention{rest} is an analytical fork. Under a quantity-noun analysis, its determiner and agreement behaviour can enter the candidate-member comparison. Under a part/subset analysis, it remains a useful access-transparency case but shouldn't supply evidence for the quantity-QN discovery criteria. Keeping those questions apart protects the profile claim from a lexical classification dispute.

TFRs are a second syntactic candidate. Their parent/daughter split and verb-class restriction suggest a cluster stabilized by constructional inheritance. The current corpus check supports the attributional/evidential contrast, but the result remains conditional on Kim's constructional analysis and lacks the category-discriminating judgments needed for a completed profile. The Spanish and Tsez cases show that the access relation travels; they don't show that the English QN cluster travels with it.

The form--meaning cases in §\ref{sec:3} are boundary applications. Their relations and diagnostics may support field-specific profiles, but this paper makes no claim about a shared compound, constructional, and processing cluster.

\subsection{Stabilizers and diagnostics}\label{sec:7-3}

\term{Lexical convention} is the proposed stabilizer for the QN profile. Speakers acquire item-specific slots such as \mention{a lot of} X, \mention{plenty of} X, \mention{the rest of} X, and \mention{a number of} X, together with their agreement and complement-selection behaviour. This account explains why the class is small and why boundary members and sense splits show graded behaviour. It doesn't require Boyd-style homeostasis. The broader countability analysis in \autocite{reynolds2026countability} supplies a possible causal-cluster connection, but the QN argument here doesn't depend on it.

\term{Constructional inheritance} is the proposed stabilizer for the TFR profile. The parent and attributional daughter support different access conditions, and the relevant candidate projections follow the verb-class hierarchy. That proposal is stronger than the current evidence warrants as a completed result, but it identifies a concrete source of covariance for the joint paper to test.

Slater's \autocite*{slater-2015-natural-kindness} stable property cluster (SPC) account asks whether a cluster is stably coinstantiated enough to support the inferential uses assigned to it. The QN profile is the strongest candidate for this diagnostic because determiner, complement, and agreement behaviours covary by subgroup. The TFR profile is a weaker candidate because its verb-class evidence is still pilot-level. Khalidi's \autocite*{khalidi-2018-causal-networks} causal-network account asks a different question about ordered causal stabilizers. The QN class appears conventionally stabilized rather than causally hierarchical, so SPC is the relevant assessment here; no Khalidian result is claimed for the umbrella.

Kindhood is downstream. A profile that supports stable projection may earn a local-kind label; a profile with partial support remains a candidate; a useful descriptive category may remain non-projectible. The comparative concept doesn't settle that question in advance.

\subsection{The resulting claim}\label{sec:7-4}

\term{Access transparency} is a relational comparative concept, not an entity or a natural kind. Its schema identifies a bearer, mediator, relation, diagnostic, and concealing counterpart. The four-gate protocol distinguishes genuine mediation from feature summation, external reference, and analysis-dependent cases that should remain indeterminate.

Three levels stay distinct:

\begin{itemize}
\tightlist
\item
  \term{Comparative concept}: access transparency, the cross-linguistic methodological tool defined by the schema of §\ref{sec:1}.
\item
  \term{Descriptive category}: language-particular grammatical content, such as English number-transparent QNs, Spanish partitive subjects, and Tsez long-distance agreement.
\item
  \term{Projectibility profile}: an observable base that licenses non-trivial projection to a target, stabilized by an identifiable source, with local-kind status as a downstream verdict.
\end{itemize}

The contribution is to apply projectibility-first discipline to a relational concept without treating every filled schema as a projectible kind. Applicability comes first. Profile evidence comes second. The distinction lets the paper preserve a useful comparative concept while remaining cautious about which language-particular patterns earn further inference.

\section{Conclusion}\label{sec:8}

Linguistic transparency is best understood as a family of mediated accessibility relations rather than as a single syntactic property. The umbrella schema, \mention{$X$ is transparent with respect to property $P$ for relation $R$ when $P$ remains accessible through $X$ for the purposes of $R$}, is a comparative concept: an analyst's question, applicable to any language, that asks which configurations let an internal property remain available to a grammatical or interpretive relation. The schema doesn't claim that any language has a descriptive category called \term{transparency}. It states the question.

The conclusion developed in §\ref{sec:7} is that the umbrella schema locates \term{projectibility profiles} rather than naming a single kind. The $X$/$P$/$R$ values across the family are heterogeneous: no shared stabilising source unifies them, and the projectible content lives at the level of the individual profile, not the level of the schema. The schema's job is to make the projectibility question askable; profiles earn scientific use when they identify an observable base licensing non-trivial projection to a target, stabilised by an identifiable source.

The profiles the paper has worked through:

\begin{itemize}
\tightlist
\item
  \term{Number-transparent quantificational nouns} (§\ref{sec:2-4}). A small, relatively stable category in the descriptive system of CGEL English (qua linguistic class, qua maintained-and-projectible category in §\ref{sec:7-1}'s stronger sense, once maintenance is attributed). Class membership predicts agreement override, characteristic determiner patterns, partitive-vs-non-partitive \mention{of}-complement availability, behaviour under fused-head construction, count/non-count percolation, and premodification tendencies. Stabilised by conventional lexical inheritance. The strongest candidate to pass Slater's \autocite*{slater-2015-natural-kindness} broad stability-plus-projection diagnostic among the cases examined; the within-domain audit is open.
\item
  \mention{The form-meaning accessibility family} (§§\ref{sec:3}--\ref{sec:5}). Three within-domain projectibility profiles related by analytical theme rather than by demonstrated cross-domain projection: a construction-level profile stabilised by constructional inheritance (CxG); a compound-level profile stabilised by usage-based entrenchment (Bell and Schäfer); a derivationally-complex-word-level profile stabilised by processing architecture (Heyer and Kornishova). Each licenses within-domain projection. Whether the three compose into a cross-domain projectibility cluster on shared items, and whether Khalidi \autocite*{khalidi-2018-causal-networks}-style causal-network refinement applies to that cluster, depend on within-items cross-modal evidence the present paper doesn't provide.
\end{itemize}

Coordination is excluded as a non-instance (§\ref{sec:2-3}): nothing was going to be encapsulated. Referential transparency is acknowledged as the boundary case (§\ref{sec:4-3}) where the schema's mediator is a context rather than a form; the philosophical-semantics tradition is named without being developed.

What is left open. Korean case-stacking and scrambling (§\ref{sec:2-6}) are deferred for separate co-authored development; the comparative-concept frame licenses asking the schema's question of those patterns with $R$-values different from agreement (argument-role recovery, scope, information-structure assignment), but the language-particular content is pending. The within-items cross-modal evidence that would test whether the form-meaning family has Khalidi-style causal-hierarchical structure (corpus predictors, ratings, priming at multiple SOAs, and interpretation accuracy on a single item set) is future work. Acquisition trajectory is a plausible further extension consistent with usage-based theories but not directly grounded by the present cases.

The paper's closing claim. \term{Transparency} in linguistics names a family of mediated accessibility relations that the umbrella schema locates as candidate projectibility profiles. A transparency attribution is scientifically useful when it identifies a profile: an observable base that licenses non-trivial projection to a target, stabilised by an identifiable source. Some profiles support strong projection; others weak; some none. Slater's SPC framework and Khalidi's causal-network account are diagnostic tools for assessing how well a profile supports induction. Kindhood is the downstream verdict on profiles that earn it, not the master question of the analysis.


\clearpage
\appendix
\numberwithin{table}{section} % appendix tables number A.1, B.1, ...
\begingroup
\scriptsize
\section{Number-transparent NP empirical apparatus}\label{sec:A}

This appendix collects the COCA pilot data behind the §\ref{sec:2} claims about the number-transparent quantificational noun class. The body of the paper carries the qualitative claims; the audit trail lives here. Full machine-readable data: \texttt{paper/data/coca-pilot.md}; KWIC spot-checks: \texttt{paper/data/kwic-checks.md}.

\subsection{Methodology}\label{sec:A-1}

Queries were issued through the local Playwright wrapper at \path{tools/english-corpora/bin/ecorg.mjs} (May 2026). Each query is a surface string; counts are raw COCA frequencies, not filtered for head versus modifier use, idiom contamination, or sense disambiguation, except where noted. Two filtering issues recur:

\begin{itemize}
\item
  \term{Modifier confound}. Surface strings of the form X \mention{of} N may include cases where N is a modifier inside a larger NP rather than the head of the \mention{of}-complement (e.g., \mention{lots of day-care} rather than \mention{lots of day}). The countability paper \autocite{reynolds2026countability} handles the same coding issue for \mention{three police} and \mention{three cattle}. Where small counts fall in the predicted-counter direction, the modifier confound likely inflates them.
\item
  \term{Parse-shift}. Surface strings containing a number-transparent QN plus a verb form may include cases where the QN sits inside a relative clause modifying a different head-noun, with the visible verb agreeing with that other head. \mention{A lot of money are} is the textbook case: of 13 raw plural-agreement hits, 10 are parse-shifted (people/parents/movies/times \mention{are}; \mention{a lot of money} in a relative clause), 1 is irrealis \mention{were} (CGEL §3.5.1), 1 \mention{'s} hit is contracted \mention{has} (singular, matches prediction), and 1 is a genuine plural override.
\end{itemize}

The wrapper's KWIC mode is unreliable for API-driven retrieval from this script as of May 2026 (Playwright internal error on every attempted call), but that doesn't mean KWIC is unavailable for human use. Counts in the tables below were generated via list-mode queries; KWIC checks were then used to filter selected small cells where raw surface counts were likely to overstate the target construction.

\subsection{\texorpdfstring{A.2 Partitive vs.~non-partitive \mention{of}-complements}{A.2 Partitive vs.~non-partitive of-complements}}\label{a.2-partitive-vs.-non-partitive-of-complements}

The contrast: \term{partitive} \mention{of}-complements have definite complements (\mention{of the money}, \mention{of those books}); \term{non-partitive} \mention{of}-complements have indefinite complements (\mention{of money}, \mention{of soldiers}, \mention{of some paintings}).

\subsubsection{\texorpdfstring{Plural-only forms (\mention{lots, bags, heaps, loads, oodles, stacks, piles})}{Plural-only forms (lots, bags, heaps, loads, oodles, stacks, piles)}}\label{plural-only-forms-lots-bags-heaps-loads-oodles-stacks-piles}

Non-partitive (X of N): \mention{lots of people} 3,878; \mention{lots of money} 1,376; \mention{lots of time} 703; \mention{loads of money} 109; \mention{piles of money} 91; \mention{loads of people} 90; \mention{loads of time} 34; \mention{heaps of money} 15; \mention{heaps of people} 7; plus single-digit hits for several other combinations. \textbf{Total: 6,311.} KWIC checks of the single-digit cells show the expected mixture: \mention{lots of day} is modifier use (\mention{day-to-day}, \mention{day hikes}, \mention{day trips}); \mention{heaps of time} and \mention{piles of time} are genuine non-partitive mass-complement uses; \mention{piles of people} is literal piling rather than the quantificational use.

Partitive (X of the N): \mention{lots of the time} 7; \mention{lots of the people} 5; \mention{loads of the day} 1. \textbf{Total: 13.} \mention{Piles} and \mention{heaps} return 0 partitive hits with any complement tested.

Indefinite-NP complement (X of a N): \mention{lots of a day} 0; \mention{piles of a day} 0; \mention{heaps of a day} 0; \mention{loads of a day} 0; \mention{lots/piles/heaps/loads of an afternoon} 0. \textbf{Total: 0.} Plural-only forms reject \mention{of} [Det NP] whether the determiner is definite (\mention{the}) or indefinite (\mention{a/an}); their \mention{of}-complements are bare/non-determined.

Partitive ratio: $13 / (6{,}311 + 13) \approx\! 0.2\%$.

\subsubsection{\texorpdfstring{\mention{Plenty}}{Plenty}}\label{plenty}

Non-partitive: \mention{plenty of time} 4,017; \mention{plenty of people} 1,563; \mention{plenty of money} 668; \mention{plenty of day} 4 (likely modifier). \textbf{Total: 6,252.}

Partitive: \mention{plenty of the time} 1; \mention{plenty of the people} 1. \textbf{Total: 2.}

Partitive ratio: $2 / 6{,}254 \approx\! 0.03\%$.

\subsubsection{\texorpdfstring{\mention{Rest, remainder}}{Rest, remainder}}\label{rest-remainder}

Non-partitive (the X of N): \mention{the rest of time} 75; \mention{the rest of people} 12; \mention{the rest of day} 7; \mention{the remainder of time} 3; \mention{the remainder of day} 2; \mention{the rest of money} 1. \textbf{Total: 100 raw.} KWIC checks confirm that this raw total is conservative. \mention{The rest of time} is mostly fixed or duration-level temporal use (\mention{for the rest of time}, \mention{spend the rest of time}), not the partitive-count contrast targeted here; \mention{the rest of people} contains possessives and non-standard \enquote{other people} uses; \mention{the rest of day} is temporal ellipsis or \mention{day} as a modifier; and the \mention{remainder} hits are technical or duration expressions.

Partitive (the X of the N): \mention{the rest of the day} 1,932; \mention{the rest of the time} 830; \mention{the rest of the people} 272; \mention{the rest of the money} 222; \mention{the remainder of the day} 75; \mention{the remainder of the time} 21; \mention{the remainder of the money} 13. \textbf{Total: 3,365.}

Indefinite-NP complement: \mention{the remainder of a day} 1.

Partitive (definite NP) ratio: $3{,}365 / (3{,}365 + 100 + 1) \approx\! 97.1\%$ raw, higher after KWIC filtering of the non-target temporal, possessive, and generic uses in the raw non-partitive cells.

\subsection{Override-direction agreement}\label{sec:A-3}

Test of the \enquote{near-obligatory in central cases} claim. Subject NP plus finite verb agreement, by complement number/count.

{\def\LTcaptype{none} % don't increment counter
\begin{longtable}[]{@{}
  >{\raggedright\arraybackslash}p{(\linewidth - 8\tabcolsep) * \real{0.1667}}
  >{\raggedright\arraybackslash}p{(\linewidth - 8\tabcolsep) * \real{0.1667}}
  >{\raggedleft\arraybackslash}p{(\linewidth - 8\tabcolsep) * \real{0.2222}}
  >{\raggedleft\arraybackslash}p{(\linewidth - 8\tabcolsep) * \real{0.2222}}
  >{\raggedleft\arraybackslash}p{(\linewidth - 8\tabcolsep) * \real{0.2222}}@{}}
\toprule\noalign{}
\begin{minipage}[b]{\linewidth}\raggedright
Subject
\end{minipage} & \begin{minipage}[b]{\linewidth}\raggedright
Predicted
\end{minipage} & \begin{minipage}[b]{\linewidth}\raggedleft
Plural agr.
\end{minipage} & \begin{minipage}[b]{\linewidth}\raggedleft
Singular agr.
\end{minipage} & \begin{minipage}[b]{\linewidth}\raggedleft
\% predicted
\end{minipage} \\
\midrule\noalign{}
\endhead
\bottomrule\noalign{}
\endlastfoot
\mention{a lot of people} & plural & 4,195 (ARE 3,406; WERE 789) & 85 (IS 72; WAS 13) & 98.0\% \\
\mention{a lot of money} & singular & 1 (ARE, filtered) & 90 (IS 61; WAS 29) & 98.9\% \\
\mention{a number of people} & plural & 100 (ARE 52; WERE 48) & 0 (4 false positives) & 100\% \\
\mention{lots of people} & plural & 348 (ARE 288; WERE 60) & 0 (6 false positives/non-standard) & 100\% \\
\mention{lots of money} & singular & 0 (1 false positive) & 24 (IS 19; WAS 5) & 100\% \\
\mention{plenty of people} & plural & 79 (ARE 65; WERE 14) & 0 & 100\% \\
\mention{plenty of money} & singular & 0 & 6 (IS 4; WAS 2) & 100\% \\
\mention{the rest of the people} & plural & 19 (ARE 15; WERE 4) & 0 & 100\% \\
\mention{the rest of the money} & singular & 0 & 19 (IS 10; WAS 9) & 100\% \\
\end{longtable}
}

Override is in the predicted direction at 98--100\% across all tested QN+complement combinations after filtering. \mention{A lot of money} required filtering: of 13 raw plural-agreement hits, 10 were parse-shifted (the surface string \mention{a lot of money are} appearing inside relative clauses modifying a plural-noun head, e.g., \mention{people who earn a lot of money are successful}; \mention{movies that make a lot of money are the biggest help}; \mention{times when we\ldots{} make a lot of money are gone}), 1 was the irrealis \mention{were} (\mention{if a lot of money were at stake}), and 1 \mention{'s} hit was singular \mention{has} (matches prediction); only 1 hit was a genuine plural override (\mention{a lot of money are being raised}). Filtered count: 1 plural / 90 singular = 98.9\% singular.

The other checked counter-direction cells are artefacts. \mention{A number of people is/was} is headed by \mention{condition}, \mention{suggestion}, \mention{ability}, or \mention{counsel}, not by \mention{number}; \mention{lots of people is/was} is headed by expressions such as \mention{murder}, \mention{hauling}, or \mention{disappearance}, with one non-standard existential hit; and \mention{lots of money are} occurs inside a PP modifying plural \mention{men}. The large \mention{a lot of people is/was} cell remains unfiltered, so the 98.0\% row is conservative.

\subsection{\texorpdfstring{A.4 \mention{Bunch}: animate vs.~inanimate aggregates}{A.4 Bunch: animate vs.~inanimate aggregates}}\label{a.4-bunch-animate-vs.-inanimate-aggregates}

CGEL claims that \mention{bunch} tilts collective for inanimate aggregates and transparent for groups of people. The COCA pilot supports the animate half but not the inanimate half.

{\def\LTcaptype{none} % don't increment counter
\begin{longtable}[]{@{}lrrr@{}}
\toprule\noalign{}
Subject & Total tokens & \mention{was} & \mention{were} \\
\midrule\noalign{}
\endhead
\bottomrule\noalign{}
\endlastfoot
\textbf{Animate aggregates} & & & \\
\mention{a bunch of people} & 1,421 & 0 & 28 \\
\mention{a bunch of kids} & 448 & 0 & 12 \\
\mention{a bunch of hooligans} & 12 & 0 & 0 \\
\textbf{Inanimate aggregates} & & & \\
\mention{a bunch of flowers} & 68 & 0 & 0 \\
\mention{a bunch of cars} & -- & 0 & 1 \\
\mention{a bunch of papers} & -- & 0 & 0 \\
\mention{a bunch of leaves} & -- & 0 & 0 \\
\mention{a bunch of things} & -- & 0 & 1 \\
\end{longtable}
}

Animates: 40 plural / 0 singular = 100\% plural. KWIC checks found no parse-shift confound in these animate-plural hits. The CGEL \enquote{transparent for groups of people} claim is supported.

Inanimates: 2 plural / 0 singular. The CGEL \enquote{tilts collective for inanimate aggregates} claim isn't supported in COCA. Inanimate-bunch is mostly absent from subject position (typical use is as object: \mention{she gave him a bunch of flowers}); the two inanimate-bunch hits in subject position both took plural rather than the CGEL-predicted singular. \mention{A bunch of flowers was presented to the teacher} (the CGEL example) returns 0 hits in COCA.

\endgroup
\begingroup
\footnotesize
\section{Accessibility topologies beyond the worked cases}\label{sec:B}

Spanish and Tsez (§\ref{sec:6-topologies}) are the cross-linguistic cases the body of the paper works. This appendix sets out four further accessibility topologies the schema maps, each with the applicability or status condition that would weaken or retire it, and states their evidential scope. None is promoted to a worked projectibility profile; the point is to show the schema's reach and its edges, and to mark which entries lack a concealing control. Every case remains subject to the §\ref{sec:2-3} counterpart test: identify $X$, $P$, and $R$, then ask what same-language pattern would make the proposed access disappear.

{
\begin{longtable}[]{@{}
  >{\raggedright\arraybackslash}p{(\linewidth - 4\tabcolsep) * \real{0.23}}
  >{\raggedright\arraybackslash}p{(\linewidth - 4\tabcolsep) * \real{0.42}}
  >{\raggedright\arraybackslash}p{(\linewidth - 4\tabcolsep) * \real{0.35}}@{}}
\caption{Accessibility topologies under the umbrella schema.}\label{tab:topologies}\\
\toprule\noalign{}
Topology & $X$--$P$--$R$ binding & Applicability condition and status \\
\midrule\noalign{}
\endfirsthead
\toprule\noalign{}
Topology & $X$--$P$--$R$ binding & Applicability condition and status \\
\midrule\noalign{}
\endhead
\bottomrule\noalign{}
\endlastfoot
Outward endpoint access & $X$ = NP or clause mediator; $P$ = complement or embedded feature; $R$ = external agreement. English QNs and transparent free relatives, Spanish partitives, and Tsez long-distance agreement instantiate this topology. & Fails where only the mediator's own head/class feature controls $R$, or where the construction-type has no concealing variant. The paper's core positive pattern. \\
\addlinespace
Null-source subpart access & $X$ = predicate configuration containing a possessed NP with an unpronounced possessor; $P$ = the recoverable possessor's person/number; $R$ = anaphoric verbal morphology plus subject case. Aleut is the key case \autocite[195]{merchant-2011-aleut-case-matters}. & Fails if the omitted possessor patterns like an ordinary overt possessor, or if the anaphoric verbal suffix appears without the missing subpart. Positive with a qualification: Merchant analyzes the missing possessor as syntactically present pro. \\
\addlinespace
Inward dependent access & $X$ = clause or nominal-dependency domain; $P$ = clausal TAM/case value; $R$ = morphology on nominal dependents. Lardil and Kayardild show clausal TAM values on dependent nominals \autocite{nordlinger-sadler-2004-tense-beyond-verb}. & Fails if dependent marking is only local nominal morphology, with no covariation with clausal predication. Candidate topology; application indeterminate until an independently motivated same-family concealing counterpart is established. \\
\addlinespace
Path-distributed access & $X$ = extraction path; $P$ = grammatical relation/case of the extracted item; $R$ = extraction-agreement morphology along that path. Chamorro supplies the comparison \autocite{chung-1998-design-agreement,chung-2004-restructuring-chamorro}. & Fails if the morphology appears only at the gap or endpoint, or if it doesn't vary with the extracted item's role. Candidate topology; application indeterminate until an independently motivated same-family concealing counterpart is established. \\
\addlinespace
Feature relay boundary case & $X$ = possessed NP; $P$ = possessor obviation/proximation; $R$ = possessed-noun morphology and verbal obviation agreement. Kutenai and Algonquian obviation are useful near misses \autocite{dryer-1992-kutenai-algonquian-obviation}. & Fails as transparency where the predicate sees only a feature already relayed onto the possessed head. Boundary case: transformed access, outside the direct-mediated pattern. \\
\end{longtable}
}

The table is a topology map, not a second set of worked cases. Aleut is the qualified positive; the head-marked and obviation entries remain conditional on the analyses and status conditions stated in the table. Lardil, Kayardild, and Chamorro mark candidate applications rather than established instances because an independently motivated same-family concealing control hasn't been shown.


\section*{Data availability}

The QN corpus material behind §\ref{sec:2} and Appendix~\ref{sec:A-3} is the COCA and COHA pilot described there: query specifications, aggregate counts, raw KWIC checks, coding notes, and the pre-registered protocol and results for the \mention{piles} candidate-member test. These files travel with the paper as supplementary material and are versioned in the project repository at \url{https://github.com/BrettRey/Linguistic_transparency}. The supplementary README maps the reported QN counts to their queries, and the repository version provides the reproducibility anchor for the current draft.

The TFR material behind §\ref{sec:2-5} (query archives, row-level coding, and the 22 August 2026 verification audit) is archived separately with the companion working paper \autocite{kim-reynolds-2026-two-kinds-transparency}.

\endgroup
\clearpage
\section*{Declaration of generative AI and AI-assisted technologies in the manuscript preparation process}

During preparation of this manuscript and the corpus work reported in §\ref{sec:2-5}, the authors used Claude Opus 4.6, Claude Opus 4.7, Claude Opus 5 (1M context), Claude Fable 5, Ox Alpha, OpenAI Codex 5.5, OpenAI GPT-5.4, and GPT-5.6 Sol (pro) in OpenAI Codex for drafting, editing, review, and first-pass corpus classification. A second model-assisted pass checked the named TFR tokens, all PP cases, a fixed-seed census sample, and the re-archived March exclusion queries against the corpus archive. After using these tools, the authors reviewed and edited the outputs as needed and take full responsibility for the content.

\begingroup
\scriptsize
\printbibliography
\endgroup

\end{document}
